\chapter{ความเป็นทวิภาคคลื่น-อนุภาคและการแทรกสอดควอนตัม}
\label{chap:wave_particle_duality}
\index{ทวิภาคคลื่น-อนุภาค (Wave-Particle Duality)}
\index{คลื่นสสารเดอ บรอยล์ (de Broglie Matter Wave)}
\index{หลักความไม่แน่นอนของไฮเซนเบิร์ก (Heisenberg Uncertainty Principle)}
\index{การแทรกสอดช่องคู่ (Double-Slit Interference)}
\index{การเลี้ยวเบนของอิเล็กตรอน (Electron Diffraction)}
\index{กลุ่มคลื่นเกาส์เซียน (Gaussian Wave Packet)}
\index{อินเทอร์เฟอโรเมทรีของโมเลกุลขนาดใหญ่ (Macromolecule Interferometry)}
\index{การทดลองลบควอนตัม (Quantum Eraser)}

% ==========================================================
% แผนบริหารการสอนประจำบทที่ 2
% ==========================================================
\section*{แผนบริหารการสอนประจำบทที่ 2}
\addcontentsline{toc}{section}{แผนบริหารการสอนประจำบทที่ 2}

\begin{lessonplanbox}
\noindent\textbf{หัวข้อเนื้อหาประจำบท}
\begin{enumerate}
    \item ภูมิหลังทางแนวคิดและสมมติฐานคลื่นสสารของหลุยส์ เดอ บรอยล์
    \item การอนุพัทธ์ความยาวคลื่นเดอ บรอยล์ในกรณีคลาสสิกและสัมพัทธภาพ
    \item การยืนยันเชิงประจักษ์ ---- การทดลองเดวิสสัน-เจอร์เมอร์ และการทดลองของจี.พี. ทอมสัน
    \item จลนศาสตร์ของคลื่นสสาร ความเร็วเฟส ความเร็วกลุ่ม และการกระจายตัวของกลุ่มคลื่น
    \item การอนุพัทธ์กลุ่มคลื่นเกาส์เซียนและการขยายตัวเชิงเวลาของความกว้างคลื่น
    \item การอนุพัทธ์หลักความไม่แน่นอนของไฮเซนเบิร์กจากความไม่เท่ากันของชวาร์ซ
    \item การแทรกสอดช่องคู่ของอนุภาคเดี่ยว การยุบตัวของฟังก์ชันคลื่น และการทดลองของโทโนมูระ
    \item พรมแดนการทดลองยุคใหม่ ---- อินเทอร์เฟอโรเมทรีของโมเลกุลชีวภาพขนาดใหญ่ การทดลองเลือกแบบหน่วงเวลาของวีลเลอร์ และการลบควอนตัม
    \item การประยุกต์เทคโนโลยี ---- กล้องจุลทรรศน์อิเล็กตรอนแบบส่องผ่าน (TEM) และกล้องไครโอ-อีเอ็ม
    \item ห้องปฏิบัติการเสมือนจริง WebAR การจำลองเชิงตัวเลขด้วย Python และการประยุกต์ P-S-C-T
\end{enumerate}

\noindent\textbf{วัตถุประสงค์เชิงพฤติกรรม}
\begin{enumerate}
    \item อธิบายสมมติฐานคลื่นสสารของเดอ บรอยล์ และคำนวณความยาวคลื่นของอนุภาคทั้งในกรณีพลังงานต่ำและกรณีสัมพัทธภาพได้อย่างถูกต้อง
    \item วิเคราะห์เงื่อนไขการเลี้ยวเบนของแบรกก์ในการทดลองเดวิสสัน-เจอร์เมอร์ และเปรียบเทียบผลลัพธ์เชิงทฤษฎีกับการทดลองได้อย่างแม่นยำ
    \item พิสูจน์ความสัมพันธ์ระหว่างความเร็วกลุ่มของคลื่นสสารกับความเร็วของอนุภาค ($v_g = v$) และอธิบายว่าเหตุใดความเร็วเฟสจึงมากกว่าความเร็วแสงได้โดยไม่ละเมิดสัมพัทธภาพ
    \item อนุพัทธ์อสมการความไม่แน่นอนของไฮเซนเบิร์ก $\Delta x \Delta p \ge \hbar/2$ จากพีชคณิตตัวดำเนินการและความไม่เท่ากันของโคชี-ชวาร์ซได้อย่างรัดกุม
    \item อธิบายกลไกการแทรกสอดของโมเลกุลชีวภาพขนาดใหญ่ (Gramicidin) และผลของการสูญเสียความพร้อมเพรียงเชิงควอนตัม (Quantum Decoherence) ผ่านการแผ่รังสีความร้อนได้อย่างลึกซึ้ง
\end{enumerate}
\end{lessonplanbox}

% ==========================================================
% 2.1 สมมติฐานคลื่นสสารของเดอ บรอยล์
% ==========================================================
\newpage
\section{สมมติฐานคลื่นสสารของเดอ บรอยล์}
\index{คลื่นสสารเดอ บรอยล์ (de Broglie Matter Wave)}
\index{ความยาวคลื่นเดอ บรอยล์ (de Broglie Wavelength)}
\index{ความสมมาตรของธรรมชาติ (Symmetry of Nature)}

\subsection{ภูมิหลังและแรงบันดาลใจจากความสมมาตรของเอกภพ}

ภายหลังการค้นพบของพลังค์และไอน์สไตน์ที่ยืนยันว่า คลื่นแม่เหล็กไฟฟ้าซึ่งแต่เดิมถือเป็นคลื่นต่อเนื่องตามทฤษฎีของแมกซ์เวลล์ มีพฤติกรรมเป็นอนุภาคพลังงาน (โฟตอน) ได้ คำถามสำคัญระดับกระบวนทัศน์จึงเกิดขึ้นว่า \textbf{ความสมมาตรของธรรมชาตินี้เกิดขึ้นในทิศทางตรงกันข้ามได้หรือไม่}

ในคริสต์ศักราช 1924 เจ้าชายหลุยส์ วิกตอร์ เดอ บรอยล์ (Louis-Victor-Pierre-Raymond, 7th duc de Broglie, 1892--1987) นักฟิสิกส์เชิงทฤษฎีชาวฝรั่งเศส ได้นำเสนอวิทยานิพนธ์ระดับปริญญาเอกอันลือลั่นต่อคณะวิทยาศาสตร์ มหาวิทยาลัยปารีส (ซอร์บอนน์) โดยเสนอสมมติฐานว่า

\begin{quote}
\textit{หากการแผ่รังสีแม่เหล็กไฟฟ้าซึ่งมีสมบัติของคลื่นสามารถแสดงพฤติกรรมเป็นอนุภาคได้ สสารทั้งหลายซึ่งเคยถูกถือว่าเป็นอนุภาค ก็ย่อมต้องมีคลื่นทางกายภาพผูกพันอยู่กับการเคลื่อนที่ของมันด้วยเช่นกัน}
\end{quote}

ความสัมพันธ์ระหว่างสมบัติเชิงอนุภาค (พลังงาน $E$ และโมเมนตัม $p$) กับสมบัติเชิงคลื่น (ความถี่ $\nu$ และความยาวคลื่น $\lambda$) ของโฟตอนคือ $E = h\nu$ และ $p = h/\lambda$ เดอ บรอยล์ได้ตั้งสมมติฐานว่าความสัมพันธ์เดียวกันนี้จะต้องใช้ได้กับอนุภาคของสสารทุกชนิด โดยที่คลื่นที่ผูกพันกับอนุภาคเรียกว่า \textbf{คลื่นสสาร (Matter Waves)} หรือ \textbf{คลื่นเดอ บรอยล์ (de Broglie Waves)} ซึ่งมีความยาวคลื่นเท่ากับ

\begin{equation}
\lambda = \frac{h}{p}
\label{eq:de_broglie_fundamental}
\end{equation}

และมีความถี่เชิงคลื่นเท่ากับ

\begin{equation}
\nu = \frac{E}{h} \implies \omega = \frac{E}{\hbar}
\label{eq:de_broglie_frequency}
\end{equation}

โดยที่ $E$ คือพลังงานรวมเชิงสัมพัทธภาพของอนุภาค ($E = \sqrt{p^2 c^2 + m^2 c^4}$)

\subsection{การคำนวณความยาวคลื่นเดอ บรอยล์ในกรณีพลังงานต่ำและสัมพัทธภาพ}

\textbf{กรณีที่ 1 อนุภาคที่เคลื่อนที่ด้วยความเร็วต่ำกว่าแสงมาก ($v \ll c$)}

ในขอบเขตกลศาสตร์ดั้งเดิม โมเมนตัมของอนุภาคคือ $p = m v$ และพลังงานจลน์คือ $K = \frac{p^2}{2m}$ ดังนั้น

\begin{equation}
\lambda = \frac{h}{mv} = \frac{h}{\sqrt{2mK}}
\label{eq:de_broglie_classical_K}
\end{equation}

สำหรับอนุภาคที่มีประจุ $q$ (เช่น อิเล็กตรอน $q = e$) ที่ถูกเร่งจากภาวะหยุดนิ่งผ่านความต่างศักย์ไฟฟ้า $V$ พลังงานจลน์คือ $K = q V$

\begin{equation}
\lambda = \frac{h}{\sqrt{2m q V}}
\label{eq:de_broglie_accelerated}
\end{equation}

แทนค่าคงตัวมูลฐานสำหรับอิเล็กตรอน ($m_e = 9.10938 \times 10^{-31}~\text{kg}$, $e = 1.60218 \times 10^{-19}~\text{C}$, $h = 6.62607 \times 10^{-34}~\text{J}\cdot\text{s}$) จะได้สูตรคำนวณสะดวกรวดเร็ว

\begin{equation}
\lambda_e = \sqrt{\frac{h^2}{2 m_e e V}} = \sqrt{\frac{1.505}{V}}~\text{nm} = \frac{1.226}{\sqrt{V}}~\text{nm}
\label{eq:de_broglie_electron_shortcut}
\end{equation}

ตัวอย่างเช่น หากเร่งอิเล็กตรอนผ่านความต่างศักย์ $V = 100~\text{V}$ จะได้ความยาวคลื่นเดอ บรอยล์ $\lambda = \frac{1.226}{\sqrt{100}} = 0.1226~\text{nm} = 1.226~\text{\AA}$ ซึ่งมีขนาดใกล้เคียงกับระยะห่างระหว่างระนาบอะตอมในผลึกของแข็ง จึงสามารถเกิดการเลี้ยวเบนได้เช่นเดียวกับรังสีเอกซ์

\textbf{กรณีที่ 2 อนุภาคที่เคลื่อนที่ด้วยความเร็วสัมพัทธภาพ ($v \sim c$)}

เมื่อพลังงานจลน์ของอนุภาคเริ่มใกล้เคียงหรือมากกว่าพลังงานนิ่ง ($K \gtrsim m c^2$) จำเป็นต้องใช้ความสัมพันธ์สัมพัทธภาพพิเศษ พลังงานรวมคือ $E = K + m c^2 = \sqrt{p^2 c^2 + m^2 c^4}$ ย้ายข้างและจัดรูปหาโมเมนตัม $p$

\begin{equation}
p^2 c^2 = E^2 - m^2 c^4 = (K + m c^2)^2 - m^2 c^4 = K^2 + 2 m c^2 K
\end{equation}

\begin{equation}
p = \frac{\sqrt{K(K + 2 m c^2)}}{c}
\label{eq:relativistic_momentum_K}
\end{equation}

ดังนั้น ความยาวคลื่นเดอ บรอยล์เชิงสัมพัทธภาพที่แท้จริงคือ

\begin{equation}
\lambda = \frac{hc}{\sqrt{K(K + 2 m c^2)}} = \frac{h}{\sqrt{2mK}}\left( 1 + \frac{K}{2mc^2} \right)^{-1/2}
\label{eq:de_broglie_relativistic_exact}
\end{equation}

ในกล้องจุลทรรศน์อิเล็กตรอนแบบส่องผ่าน (TEM) ที่ใช้ความต่างศักย์เร่งสูง $V = 200~\text{kV}$ ถึง $300~\text{kV}$ พลังงานจลน์ของอิเล็กตรอน $K = 200~\text{keV}$ คิดเป็นประมาณ $40\%$ ของพลังงานนิ่ง ($m_e c^2 = 511~\text{keV}$) หากใช้สูตรคลาสสิกจะให้ค่า $\lambda = 0.00274~\text{nm}$ แต่หากใช้สูตรสัมพัทธภาพที่ถูกต้องสมการ~\eqref{eq:de_broglie_relativistic_exact} จะได้ $\lambda = 0.00251~\text{nm}$ ซึ่งแตกต่างกันถึง $9.2\%$ อันมีผลอย่างมีนัยสำคัญต่อการคำนวณการเลี้ยวเบนของอิเล็กตรอนในโครงสร้างผลึก

% ==========================================================
% 2.2 การยืนยันเชิงประจักษ์
% ==========================================================
\section{การยืนยันเชิงประจักษ์ของการเลี้ยวเบนคลื่นสสาร}
\index{การทดลองเดวิสสัน-เจอร์เมอร์ (Davisson-Germer Experiment)}
\index{การเลี้ยวเบนของแบรกก์ (Bragg Diffraction)}
\index{การทดลองของจี.พี. ทอมสัน (G.P. Thomson Experiment)}

\subsection{การทดลองเดวิสสัน-เจอร์เมอร์และผลึกนิกเกิล}

การพิสูจน์สมมติฐานของเดอ บรอยล์เกิดขึ้นอย่างเด็ดขาดในคริสต์ศักราช 1927 โดย คลินตัน เดวิสสัน (Clinton Davisson) และ เลสเตอร์ เจอร์เมอร์ (Lester Germer) ณ ห้องปฏิบัติการเบลล์ (Bell Telephone Laboratories) ประเทศสหรัฐอเมริกา

ในการทดลอง ลำอิเล็กตรอนพลังงานต่ำถูกยิงเข้าใส่เป้าผลึกเดี่ยวของนิกเกิล (Single Crystal Nickel) ดังแสดงในภาพที่~\ref{fig:davisson_germer_apparatus} เมื่อปรับเปลี่ยนพลังงานของอิเล็กตรอนและมุมตรวจวัด ($\phi$) พบว่าที่พลังงานจลน์ $K = 54~\text{eV}$ (ศักย์เร่ง $54~\text{V}$) มีลำอิเล็กตรอนกระเจิงออกมาอย่างเข้มข้นเด่นชัดที่มุม $\phi = 50.0^\circ$

\begin{figure}[H]
\centering
\begin{tikzpicture}[scale=1.1]
    % Nickel crystal block
    \draw[line width=1.5pt, darkNavy, fill=slateGrey!25] (0, -0.5) rectangle (5.0, 0.5);
    \node at (2.5, 0) [font=\small\bfseries\color{darkNavy}] {ผลึกเดี่ยวนิกเกิล (Nickel Crystal)};

    % Atomic layers
    \foreach \x in {0.8, 1.4, 2.0, 2.6, 3.2, 3.8, 4.4} {
        \fill[quantumGold] (\x, 0.5) circle (3pt);
        \fill[quantumGold!70] (\x, -0.1) circle (3pt);
    }
    \node at (5.8, 0.5) [font=\scriptsize\color{quantumGold!90!black}] {ระนาบอะตอม};

    % Electron Gun
    \draw[line width=1.8pt, darkNavy, fill=cyberBlue!10] (-2.0, 3.5) rectangle (-0.8, 4.3);
    \node at (-1.4, 4.6) [font=\scriptsize\bfseries\color{darkNavy}] {ปืนยิงอิเล็กตรอน};
    \draw[-{Stealth[scale=1.3]}, line width=2pt, cyberBlue] (-1.4, 3.5) -- (2.6, 0.5);
    \node at (0.2, 2.5) [font=\small\bfseries\color{cyberBlue}] {ลำอิเล็กตรอน ($54~\text{eV}$)};

    % Scattered Electron Beam
    \draw[-{Stealth[scale=1.3]}, line width=2pt, entanglePink] (2.6, 0.5) -- (4.8, 3.2);
    \node at (4.5, 1.8) [font=\small\bfseries\color{entanglePink}] {ลำสะท้อนเข้มข้น};

    % Movable Detector
    \draw[line width=1.5pt, neonPurple, fill=neonPurple!15] (4.6, 3.0) circle (0.4cm);
    \node at (5.7, 3.5) [font=\scriptsize\bfseries\color{neonPurple}] {ตัวตรวจจับปรับมุมได้};

    % Scattering angle arc
    \draw[dashed, slateGrey, line width=1pt] (2.6, 0.5) -- (2.6, 3.5);
    \draw[-{Stealth}, line width=1.2pt, neonPurple] (2.6, 2.0) arc (90:50:1.5);
    \node at (3.2, 2.3) [font=\small\bfseries\color{neonPurple}] {$\phi = 50^\circ$};
\end{tikzpicture}
\caption{แผนผังการทดลองเดวิสสัน-เจอร์เมอร์ แสดงลำอิเล็กตรอนตกกระทบตั้งฉากกับผลึกนิกเกิลและเกิดการกระเจิงแบบแทรกสอดเสริมกันที่มุม 50 องศา}
\label{fig:davisson_germer_apparatus}
\end{figure}

จากข้อมูลโครงสร้างผลึกด้วยรังสีเอกซ์ ระยะห่างระหว่างแถวอะตอมบนผิวผลึกนิกเกิลคือ $D = 0.215~\text{nm}$ เงื่อนไขการแทรกสอดเสริมกันเชิงสร้างสรรค์ (Constructive Interference) สำหรับการสะท้อนจากตะแกรงผิวคือ

\begin{equation}
n \lambda = D \sin\phi
\label{eq:surface_diffraction}
\end{equation}

สำหรับอันดับที่หนึ่ง ($n = 1$) จะได้ความยาวคลื่นจากการทดลองคือ

\begin{equation}
\lambda_{\text{exp}} = (0.215~\text{nm})\sin(50.0^\circ) = (0.215)(0.7660) = 0.165~\text{nm}
\end{equation}

ในขณะเดียวกัน ความยาวคลื่นเดอ บรอยล์ที่คำนวณตามทฤษฎีจากสมการ~\eqref{eq:de_broglie_electron_shortcut} สำหรับ $V = 54~\text{V}$ คือ

\begin{equation}
\lambda_{\text{theory}} = \frac{1.226}{\sqrt{54}}~\text{nm} = \frac{1.226}{7.348}~\text{nm} = 0.1668~\text{nm} \approx 0.167~\text{nm}
\end{equation}

ความสอดคล้องอย่างยอดเยี่ยมระหว่างผลการทดลอง ($0.165~\text{nm}$) และทฤษฎีคลื่นสสาร ($0.167~\text{nm}$) ด้วยความคลาดเคลื่อนน้อยกว่า $1\%$ ถือเป็นหลักฐานเชิงประจักษ์แรกที่ยืนยันว่าอนุภาคสสารมีสมบัติเป็นคลื่นอย่างแท้จริง

\subsection{การทดลองของจี.พี. ทอมสันและลวดลายการเลี้ยวเบนทรงกลม}

ในปีเดียวกัน คริสต์ศักราช 1927 จอร์จ แพเก็ต ทอมสัน (George Paget Thomson, 1892--1975) บุตรชายของ เจ.เจ. ทอมสัน (ผู้ค้นพบอนุภาคอิเล็กตรอน) ได้ทำการทดลองยิงลำอิเล็กตรอนพลังงานสูง ($10--60~\text{keV}$) ทะลุผ่านแผ่นฟอยล์ทองคำบางเฉียบ (Thin Polycrystalline Gold Foil) ผลปรากฏบนแผ่นฟิล์มด้านหลังเป็น \textbf{วงแหวนการเลี้ยวเบนศูนย์กลางร่วม (Concentric Diffraction Rings)} ที่สอดคล้องกับรูปแบบการเลี้ยวเบนรังสีเอกซ์ของเดอบาย-เชอร์เรอร์ (Debye-Scherrer X-ray rings) ทุกประการ

ประวัติศาสตร์ฟิสิกส์จึงจารึกไว้อย่างน่าอัศจรรย์ว่า เจ.เจ. ทอมสัน ได้รับรางวัลโนเบลในปี 1906 จากการพิสูจน์ว่า \textbf{อิเล็กตรอนเป็นอนุภาค} ในขณะที่ จี.พี. ทอมสัน บุตรชาย ได้รับรางวัลโนเบลร่วมกับเดวิสสันในปี 1937 จากการพิสูจน์ว่า \textbf{อิเล็กตรอนเป็นคลื่น} ซึ่งตอกย้ำแนวคิดทวิภาคคลื่น-อนุภาคอย่างลึกซึ้ง

% ==========================================================
% 2.3 จลนศาสตร์ของคลื่นสสาร ความเร็วเฟส และความเร็วกลุ่ม
% ==========================================================
\section{จลนศาสตร์ของคลื่นสสาร ความเร็วเฟส และความเร็วกลุ่ม}
\index{ความเร็วเฟส (Phase Velocity)}
\index{ความเร็วกลุ่ม (Group Velocity)}
\index{กลุ่มคลื่น (Wave Packet)}

\subsection{ความเร็วเฟสและปมปริศนาความเร็วเหนือแสง}

พิจารณาคลื่นสสารแบบระนาบความถี่เดี่ยว (Monochromatic Plane Wave) ในหนึ่งมิติ

\begin{equation}
\psi(x, t) = A e^{i(k x - \omega t)}
\label{eq:plane_matter_wave}
\end{equation}

โดยที่เลขคลื่น $k = \frac{2\pi}{\lambda} = \frac{p}{\hbar}$ และความถี่เชิงมุม $\omega = 2\pi\nu = \frac{E}{\hbar}$ \textbf{ความเร็วเฟส (Phase Velocity, $v_p$)} คือความเร็วในการเคลื่อนที่ของจุดที่มีเฟสคงที่ ($kx - \omega t = \text{const}$)

\begin{equation}
v_p = \frac{\omega}{k} = \frac{E/\hbar}{p/\hbar} = \frac{E}{p}
\label{eq:phase_velocity_def}
\end{equation}

ใช้ความสัมพันธ์เชิงสัมพัทธภาพ $E = \gamma m c^2$ และ $p = \gamma m v$ เมื่อ $\gamma = 1/\sqrt{1 - v^2/c^2}$ และ $v$ คือความเร็วของอนุภาค

\begin{equation}
v_p = \frac{\gamma m c^2}{\gamma m v} = \frac{c^2}{v}
\label{eq:phase_velocity_superluminal}
\end{equation}

เนื่องจากอนุภาคที่มีมวลต้องเคลื่อนที่ด้วยความเร็วช้ากว่าแสงในสุญญากาศเสมอ ($v < c$) สมการ~\eqref{eq:phase_velocity_superluminal} นำไปสู่ผลลัพธ์ที่ชวนฉงนว่า

\begin{equation}
v_p = \frac{c^2}{v} > c
\end{equation}

ความเร็วเฟสของคลื่นสสารเดอ บรอยล์มีค่ามากกว่าความเร็วแสงเสมอ อย่างไรก็ตาม สิ่งนี้ไม่ได้ละเมิดทฤษฎีสัมพัทธภาพพิเศษ เนื่องจากคลื่นระนาบความถี่เดี่ยวแบบบริสุทธิ์ทอดยาวไปทั่วทั้งปริภูมิอย่างไม่สิ้นสุด จึงไม่สามารถใช้ส่งสัญญาณสารสนเทศหรือระบุตำแหน่งของอนุภาคได้ การส่งผ่านสารสนเทศและอนุภาคที่แท้จริงต้องอาศัยการรวมกลุ่มของคลื่นหลายความถี่เข้าด้วยกันเป็น \textbf{กลุ่มคลื่น (Wave Packet)}

\subsection{ความเร็วกลุ่มและความเท่ากันกับความเร็วอนุภาค}

เมื่อรวมคลื่นระนาบที่มีความถี่และเลขคลื่นแตกต่างกันเล็กน้อย การแทรกสอดจะสร้างโครงสร้างเปลือกหุ้ม (Envelope) ที่เคลื่อนที่ด้วย \textbf{ความเร็วกลุ่ม (Group Velocity, $v_g$)}

\begin{equation}
v_g = \frac{d\omega}{dk}
\label{eq:group_velocity_def}
\end{equation}

ใช้ความสัมพันธ์ของเดอ บรอยล์ $\omega = E/\hbar$ และ $k = p/\hbar$ ดังนั้น $d\omega = dE/\hbar$ และ $dk = dp/\hbar$

\begin{equation}
v_g = \frac{dE/\hbar}{dp/\hbar} = \frac{dE}{dp}
\label{eq:group_velocity_dEdp}
\end{equation}

จากความสัมพันธ์สัมพัทธภาพ $E^2 = p^2 c^2 + m^2 c^4$ หาอนุพันธ์ย่อยเทียบกับ $p$

\begin{equation}
2 E \frac{dE}{dp} = 2 p c^2 \implies \frac{dE}{dp} = \frac{p c^2}{E}
\end{equation}

แทนค่า $p = \gamma m v$ และ $E = \gamma m c^2$

\begin{equation}
v_g = \frac{dE}{dp} = \frac{(\gamma m v) c^2}{\gamma m c^2} = v
\label{eq:group_velocity_equals_particle}
\end{equation}

ผลลัพธ์ในสมการ~\eqref{eq:group_velocity_equals_particle} งดงามอย่างยิ่งในทางฟิสิกส์ เพราะได้พิสูจน์ว่า \textbf{ความเร็วกลุ่มของกลุ่มคลื่นสสารมีค่าเท่ากับความเร็วของอนุภาคจริงทุกประการ} ($v_g = v_{\text{particle}} < c$) อนุภาคเคลื่อนที่ไปพร้อมกับกลุ่มคลื่น และผลคูณระหว่างความเร็วเฟสกับความเร็วกลุ่มมีค่าคงที่เสมอเท่ากับ

\begin{equation}
v_p \cdot v_g = \left( \frac{c^2}{v} \right) \cdot v = c^2
\label{eq:vp_vg_c2}
\end{equation}

% ==========================================================
% 2.4 การกระจายตัวของกลุ่มคลื่นเกาส์เซียน
% ==========================================================
\section{การกระจายตัวของกลุ่มคลื่นเกาส์เซียนเชิงคณิตศาสตร์}
\index{กลุ่มคลื่นเกาส์เซียน (Gaussian Wave Packet)}
\index{การกระจายตัวของคลื่น (Wave Packet Dispersion)}
\index{การแปลงฟูเรียร์ (Fourier Transform)}

\subsection{ฟังก์ชันคลื่นเกาส์เซียนเริ่มต้น}

พิจารณากลุ่มคลื่นที่อธิบายตำแหน่งของอนุภาคที่มีความไม่แน่นอนเริ่มต้น $\sigma_0$ และมีโมเมนตัมเฉลี่ย $p_0 = \hbar k_0$ ณ เวลา $t = 0$ ซึ่งมีรูปแบบฟังก์ชันคลื่นเป็น \textbf{กลุ่มคลื่นเกาส์เซียน (Gaussian Wave Packet)}

\begin{equation}
\psi(x, 0) = \left( \frac{1}{2\pi \sigma_0^2} \right)^{1/4} \exp\left( -\frac{x^2}{4\sigma_0^2} \right) e^{i k_0 x}
\label{eq:gaussian_packet_t0}
\end{equation}

ฟังก์ชันคลื่นนี้ได้รับการนอร์แมลไลซ์เรียบร้อยแล้ว ($\int_{-\infty}^\infty |\psi(x, 0)|^2 dx = 1$) และมีความหนาแน่นความน่าจะเป็นเป็นฟังก์ชันการกระจายแบบปกติที่มีค่าเฉลี่ย $\langle x \rangle = 0$ และส่วนเบี่ยงเบนมาตรฐาน $\Delta x = \sigma_0$

การกระจายในปริภูมิโมเมนตัมหาได้จากการแปลงฟูเรียร์ (Fourier Transform)

\begin{equation}
\phi(k) = \frac{1}{\sqrt{2\pi}} \int_{-\infty}^\infty \psi(x, 0) e^{-i k x} dx = \left( \frac{2\sigma_0^2}{\pi} \right)^{1/4} \exp\left[ -\sigma_0^2 (k - k_0)^2 \right]
\label{eq:momentum_distribution}
\end{equation}

ความกว้างของการกระจายในปริภูมิเลขคลื่นคือ $\Delta k = \frac{1}{2\sigma_0}$ ส่งผลให้ความไม่แน่นอนของโมเมนตัมเริ่มต้นคือ

\begin{equation}
\Delta p = \hbar \Delta k = \frac{\hbar}{2\sigma_0}
\end{equation}

เมื่อนำมาคูณกันจะได้

\begin{equation}
\Delta x \cdot \Delta p = \sigma_0 \cdot \left( \frac{\hbar}{2\sigma_0} \right) = \frac{\hbar}{2}
\label{eq:minimum_uncertainty_packet}
\end{equation}

ซึ่งหมายความว่า \textbf{กลุ่มคลื่นเกาส์เซียนคือกลุ่มคลื่นที่มีผลคูณความไม่แน่นอนต่ำสุดที่เป็นไปได้ในธรรมชาติ (Minimum Uncertainty Wave Packet)}

\subsection{วิวัฒนาการเชิงเวลาและการแผ่กว้างออก}

สำหรับอนุภาคอิสระมวล $m$ สัมพันธภาพการกระจายตัวแบบไม่สัมพัทธภาพคือ $\omega(k) = \frac{\hbar k^2}{2m}$ ฟังก์ชันคลื่น ณ เวลา $t > 0$ หาได้จากการอินทิเกรตรวมทุกองค์ประกอบฟูเรียร์

\begin{equation}
\psi(x, t) = \frac{1}{\sqrt{2\pi}} \int_{-\infty}^\infty \phi(k) e^{i(k x - \omega(k) t)} dk
\end{equation}

การคำนวณปริพันธ์เกาส์เซียนเชิงซ้อนจะได้ฟังก์ชันคลื่น ณ เวลา $t$

\begin{equation}
\psi(x, t) = \left( \frac{1}{2\pi \sigma_0^2} \right)^{1/4} \frac{1}{\sqrt{1 + i\frac{\hbar t}{2m\sigma_0^2}}} \exp\left[ -\frac{(x - v_0 t)^2}{4\sigma_0^2\left( 1 + i\frac{\hbar t}{2m\sigma_0^2} \right)} \right] \exp\left[ i(k_0 x - \omega_0 t) \right]
\label{eq:wave_packet_time_evolved}
\end{equation}

เมื่อหาขนาดกำลังสองจะได้ความหนาแน่นความน่าจะเป็น ณ เวลา $t$

\begin{equation}
P(x, t) = |\psi(x, t)|^2 = \frac{1}{\sqrt{2\pi}\sigma(t)} \exp\left[ -\frac{(x - v_0 t)^2}{2\sigma^2(t)} \right]
\label{eq:gaussian_probability_density_t}
\end{equation}

โดยที่จุดศูนย์กลางของกลุ่มคลื่นเคลื่อนที่ไปด้วยความเร็วคงที่ $v_0 = \frac{\hbar k_0}{m} = \frac{p_0}{m}$ แต่ความกว้างของกลุ่มคลื่น $\sigma(t) = \Delta x(t)$ จะขยายกว้างขึ้นเรื่อย ๆ ตามเวลาตามสมการ

\begin{equation}
\sigma(t) = \sigma_0 \sqrt{1 + \left( \frac{\hbar t}{2 m \sigma_0^2} \right)^2}
\label{eq:wave_packet_spreading}
\end{equation}

สมการ~\eqref{eq:wave_packet_spreading} แสดงปรากฏการณ์สำคัญทางควอนตัมที่เรียกว่า \textbf{การกระจายตัวของกลุ่มคลื่น (Wave Packet Dispersion)} เมื่อเวลาผ่านไปนานมาก ($t \gg \frac{2m\sigma_0^2}{\hbar}$) ความกว้างของกลุ่มคลื่นจะเพิ่มขึ้นแบบเชิงเส้นตามเวลา $\sigma(t) \approx \frac{\hbar t}{2m\sigma_0} = \frac{\Delta p}{m} t$ ซึ่งเกิดจากความจริงที่ว่ากลุ่มคลื่นประกอบด้วยคลื่นย่อยที่มีความเร็วไม่เท่ากันนั่นเอง

\subsection{กลุ่มคลื่นแอรี่และการเคลื่อนที่แบบเร่งโดยไร้การกระจายตัว}
\index{กลุ่มคลื่นแอรี่ (Airy Wave Packet)}
\index{ผลเฉลยไร้การกระจายตัว (Non-dispersing Solutions)}

เป็นเวลานานที่นักฟิสิกส์เชื่อว่ากลุ่มคลื่นของอนุภาคอิสระทุกชนิดในปริภูมิว่าง ($V = 0$) จะต้องแผ่กระจายตัวออกไปตามเวลาเสมอ ทว่าในคริสต์ศักราช 1979 ไมเคิล เบอร์รี (Sir Michael Berry) และ นันดอร์ บาลาซส์ (Nandor Balazs) ได้ค้นพบผลเฉลยอันน่าทึ่งของสมการชเรอดิงเงอร์หนึ่งมิติในรูปของ \textbf{ฟังก์ชันแอรี่ (Airy Function, $\text{Ai}$)}

\begin{equation}
\psi(x, t) = \text{Ai}\left[ \frac{B}{\hbar^{2/3}}\left( x - \frac{B^3 t^2}{4 m^2} \right) \right] \exp\left[ \frac{i B^3 t}{2 m \hbar}\left( x - \frac{B^3 t^2}{6 m^2} \right) \right]
\label{eq:airy_wave_packet}
\end{equation}

โดยที่ $B$ คือค่าคงตัวทางกายภาพ เมื่อหาความหนาแน่นความน่าจะเป็นจะได้

\begin{equation}
P(x, t) = |\psi(x, t)|^2 = \text{Ai}^2\left[ \frac{B}{\hbar^{2/3}}\left( x - \frac{B^3 t^2}{4 m^2} \right) \right]
\label{eq:airy_probability_density}
\end{equation}

สมการ~\eqref{eq:airy_probability_density} นำไปสู่ข้อสรุปอันน่าอัศจรรย์สองประการ
\begin{enumerate}
    \item \textbf{การคงรูปทรงโดยไร้การกระจายตัว (Shape-Preserving / Non-dispersing)} ---- รูปร่างของความหนาแน่นความน่าจะเป็นไม่เปลี่ยนแปลงไปตามเวลาเลย จึงไม่มีการแผ่กว้างออกเหมือนกลุ่มคลื่นเกาส์เซียน
    \item \textbf{การเคลื่อนที่แบบเร่งตัวเองในสุญญากาศ (Self-Acceleration)} ---- จุดยอดของกลุ่มคลื่นแอรี่เคลื่อนที่ตามแนววิถีพาราโบลา $x(t) = \frac{B^3 t^2}{4 m^2}$ ซึ่งมีความเร่งคงที่ $a = \frac{B^3}{2 m^2}$ แม้จะไม่มีแรงภายนอกใด ๆ มากระทำต่ออนุภาคเลย ($V = 0$)
\end{enumerate}

ในคริสต์ศักราช 2013 กลุ่มวิจัยของ อดินา โวโลค-โบลช (Adina Voloch-Bloch) ณ มหาวิทยาลัยเทลอาวีฟ ประสบความสำเร็จในการสร้างและสังเกตกลุ่มคลื่นแอรี่ของอิเล็กตรอนอิสระในกล้องจุลทรรศน์อิเล็กตรอน (ตีพิมพ์ในวารสาร \textit{Nature}) ยืนยันการเร่งตัวเองของคลื่นสสารในสุญญากาศตามการทำนายของกลศาสตร์ควอนตัมอย่างสมบูรณ์แบบ

% ==========================================================
% 2.5 การอนุพัทธ์หลักความไม่แน่นอนของไฮเซนเบิร์ก
% ==========================================================
\section{การอนุพัทธ์หลักความไม่แน่นอนของไฮเซนเบิร์ก}
\index{หลักความไม่แน่นอนของไฮเซนเบิร์ก (Heisenberg Uncertainty Principle)}
\index{ความไม่เท่ากันของชวาร์ซ (Schwarz Inequality)}
\index{ความสัมพันธ์คอมมูเทชัน (Commutation Relation)}

ในคริสต์ศักราช 1927 แวร์เนอร์ ไฮเซนเบิร์ก (Werner Karl Heisenberg, 1901--1976) ได้ค้นพบหลักการที่เป็นเสาหลักค้ำจุนกลศาสตร์ควอนตัม นั่นคือ \textbf{หลักความไม่แน่นอน (Uncertainty Principle)} ซึ่งระบุว่า เป็นไปไม่ได้ในทางกายภาพที่จะวัดตำแหน่งและโมเมนตัมของอนุภาคให้มีความแม่นยำอย่างสมบูรณ์พร้อม ๆ กัน

\subsection{การพิสูจน์ทั่วไปจากความไม่เท่ากันของโคชี-ชวาร์ซ}

พิจารณาตัวดำเนินการแฮร์มิเชียนสองตัวใด ๆ $\hat{A}$ และ $\hat{B}$ นิยามตัวดำเนินการความเบี่ยงเบน $\Delta \hat{A} = \hat{A} - \langle A \rangle$ และ $\Delta \hat{B} = \hat{B} - \langle B \rangle$ ความแปรปรวนนิยามโดย

\begin{equation}
(\Delta A)^2 = \langle (\Delta \hat{A})^2 \rangle = \langle \psi | (\Delta \hat{A})^2 | \psi \rangle
\end{equation}

นิยามสถานะเวกเตอร์สองสถานะ $|\alpha\rangle = \Delta \hat{A}|\psi\rangle$ และ $|\beta\rangle = \Delta \hat{B}|\psi\rangle$ จาก \textbf{ความไม่เท่ากันของโคชี-ชวาร์ซ (Cauchy-Schwarz Inequality)} ในปริภูมิฮิลเบิร์ต

\begin{equation}
\langle \alpha | \alpha \rangle \langle \beta | \beta \rangle \ge |\langle \alpha | \beta \rangle|^2
\label{eq:cauchy_schwarz}
\end{equation}

เนื่องจาก $\langle \alpha | \alpha \rangle = (\Delta A)^2$ และ $\langle \beta | \beta \rangle = (\Delta B)^2$ จึงได้

\begin{equation}
(\Delta A)^2 (\Delta B)^2 \ge |\langle \psi | \Delta \hat{A} \Delta \hat{B} | \psi \rangle|^2
\label{eq:variance_product}
\end{equation}

ผลคูณของตัวดำเนินการ $\Delta \hat{A} \Delta \hat{B}$ สามารถแยกออกเป็นผลรวมของส่วนสมมาตร (แอนตีคอมมิวเทเตอร์) และส่วนปฏิสมมาตร (คอมมิวเทเตอร์)

\begin{equation}
\Delta \hat{A} \Delta \hat{B} = \frac{1}{2}[\Delta \hat{A}, \Delta \hat{B}] + \frac{1}{2}\{\Delta \hat{A}, \Delta \hat{B}\}
\end{equation}

โดยที่ $[\Delta \hat{A}, \Delta \hat{B}] = [\hat{A}, \hat{B}]$ และสังเกตว่าพจน์คอมมิวเทเตอร์เป็นตัวดำเนินการแอนตีแฮร์มิเชียน (ให้ค่าคาดหวังเป็นจำนวนจินตภาพบริสุทธิ์) ในขณะที่แอนตีคอมมิวเทเตอร์เป็นแฮร์มิเชียน (ให้ค่าคาดหวังเป็นจำนวนจริง) ดังนั้น

\begin{equation}
|\langle \Delta \hat{A} \Delta \hat{B} \rangle|^2 = \left| \frac{1}{2}\langle [\hat{A}, \hat{B}] \rangle + \frac{1}{2}\langle \{\Delta \hat{A}, \Delta \hat{B}\} \rangle \right|^2 = \frac{1}{4}|\langle [\hat{A}, \hat{B}] \rangle|^2 + \frac{1}{4}|\langle \{\Delta \hat{A}, \Delta \hat{B}\} \rangle|^2
\end{equation}

เนื่องจากพจน์ที่สองยกกำลังสองมีค่ามากกว่าหรือเท่ากับศูนย์เสมอ เราจึงได้ \textbf{ความสัมพันธ์ความไม่แน่นอนของโรเบิร์ตสัน-ชเรอดิงเงอร์ (Robertson-Schr{\"o}dinger Uncertainty Relation)}

\begin{equation}
\Delta A \cdot \Delta B \ge \frac{1}{2}|\langle [\hat{A}, \hat{B}] \rangle|
\label{eq:robertson_uncertainty}
\end{equation}

\subsection{ความไม่แน่นอนของตำแหน่งและโมเมนตัม}

สำหรับตัวดำเนินการตำแหน่ง $\hat{x} = x$ และตัวดำเนินการโมเมนตัม $\hat{p}_x = -i\hbar\frac{\partial}{\partial x}$ ความสัมพันธ์คอมมูเทชันพื้นฐานคือ

\begin{equation}
[\hat{x}, \hat{p}_x] = i\hbar
\label{eq:canonical_commutation}
\end{equation}

แทนค่าลงในสมการ~\eqref{eq:robertson_uncertainty}

\begin{equation}
\Delta x \cdot \Delta p_x \ge \frac{1}{2}|i\hbar| = \frac{\hbar}{2}
\label{eq:heisenberg_position_momentum}
\end{equation}

สมการ~\eqref{eq:heisenberg_position_momentum} แสดงให้เห็นว่าความไม่แน่นอนไม่ใช่ข้อจำกัดของเทคโนโลยีเครื่องมือวัด แต่เป็น \textbf{สมบัติทางเรขาคณิตมูลฐานของธรรมชาติ} อันสืบเนื่องมาจากคุณลักษณะที่ไม่สลับที่กัน (Non-commutativity) ของตัวดำเนินการในกลศาสตร์ควอนตัม

\subsection{ความไม่แน่นอนของพลังงานและเวลา}

นอกจากความไม่แน่นอนของตำแหน่งและโมเมนตัมแล้ว ยังมีความสัมพันธ์ความไม่แน่นอนอีกรูปแบบหนึ่งที่มีความสำคัญอย่างยิ่งในฟิสิกส์นิวเคลียร์และอนุภาคพลังงานสูง นั่นคือ \textbf{ความสัมพันธ์ความไม่แน่นอนระหว่างพลังงานและเวลา (Energy-Time Uncertainty Relation)}

\begin{equation}
\Delta E \cdot \Delta t \ge \frac{\hbar}{2}
\label{eq:energy_time_uncertainty}
\end{equation}

ในที่นี้ เวลา $t$ ไม่ใช่ตัวดำเนินการในกลศาสตร์ควอนตัม แต่เป็นตัวแปรพารามิเตอร์ $\Delta t$ จึงหมายถึงช่วงเวลาอายุขัย (Lifetime, $\tau$) ของสถานะควอนตัม สมการนี้อธิบายว่า สถานะกระตุ้นของอะตอมที่มีอายุขัยสั้น $\tau$ จะมีความไม่แน่นอนของพลังงาน $\Delta E \approx \hbar/\tau$ ส่งผลให้เส้นสเปกตรัมที่แผ่ออกมามีความกว้างตามธรรมชาติ (Natural Linewidth) เสมอ และยังอนุญาตให้อนุภาคเสมือน (Virtual Particles) สามารถผุดขึ้นมาจากสุญญากาศควอนตัมได้ในช่วงเวลาสั้น ๆ $\Delta t \le \hbar/(2\Delta E)$ ตามความผันผวนของสุญญากาศ (Vacuum Fluctuations)

% ==========================================================
% 2.6 การแทรกสอดช่องคู่ของอนุภาคเดี่ยว
% ==========================================================
\section{การแทรกสอดช่องคู่ของอนุภาคเดี่ยวและการยุบตัวของฟังก์ชันคลื่น}
\index{การแทรกสอดช่องคู่ (Double-Slit Interference)}
\index{การทดลองของโทโนมูระ (Tonomura Experiment)}
\index{การยุบตัวของฟังก์ชันคลื่น (Wavefunction Collapse)}
\index{ข้อมูลเส้นทาง (Which-Way Information)}

\subsection{การทดลองในจินตนาการของไฟน์แมนสู่ความจริงในห้องปฏิบัติการ}

ริชาร์ด ไฟน์แมน (Richard P. Feynman) เคยกล่าวไว้ว่า การทดลองแทรกสอดช่องคู่ของอิเล็กตรอนบรรจุ \textbf{หัวใจสำคัญของกลศาสตร์ควอนตัม และความลึกลับประการเดียวที่แท้จริงของธรรมชาติ} ไว้ทั้งหมด

ในการทดลองช่องคู่คลาสสิกของยัง (Young's Double-Slit) แสงแสดงการแทรกสอดเกิดเป็นแถบมืดและสว่าง สำหรับอนุภาคอิเล็กตรอน หากอิเล็กตรอนเป็นเพียงลูกกระสุนปืนคลาสสิก เมื่อยิงผ่านช่องเปิดคู่สองช่อง ($S_1$ และ $S_2$) กราฟความน่าจะเป็นที่ปรากฏบนฉากควรเป็นเพียงผลบวกตรงไปตรงมาของความน่าจะเป็นที่ผ่านแต่ละช่อง

\begin{equation}
P_{\text{classical}}(x) = P_1(x) + P_2(x)
\end{equation}

ทว่าในการทดลองจริง ฟังก์ชันคลื่นของอิเล็กตรอนก่อนถึงฉากคือการซ้อนทับเชิงเส้น (Superposition) ของแอมพลิจูดความน่าจะเป็นจากทั้งสองช่องทาง $\psi(x) = \psi_1(x) + \psi_2(x)$ ความน่าจะเป็นในการตรวจพบอิเล็กตรอนบนฉากตามการแปลความหมายเชิงสถิติของมัคซ์ บอร์น (Max Born) คือ

\begin{equation}
P_{\text{quantum}}(x) = |\psi_1(x) + \psi_2(x)|^2 = |\psi_1(x)|^2 + |\psi_2(x)|^2 + 2\text{Re}[\psi_1^*(x)\psi_2(x)]
\label{eq:quantum_interference_pattern}
\end{equation}

พจน์ที่สาม $2\text{Re}[\psi_1^*(x)\psi_2(x)]$ คือ \textbf{พจน์การแทรกสอดควอนตัม (Quantum Interference Term)} ซึ่งทำให้เกิดริ้วแถบสว่างและมืดบนฉาก

ในคริสต์ศักราช 1989 อากิระ โทโนมูระ (Akira Tonomura, 1942--2012) และคณะ ณ สถาบันวิจัยฮิตาชิ ประเทศญี่ปุ่น ได้ทำการทดลองอันยอดเยี่ยมด้วยการใช้บิปริซึมอิเล็กตรอน (Electron Biprism) และยิงอิเล็กตรอนทีละตัวอย่างโดดเดี่ยว (Single-Electron Build-up) โดยเว้นระยะห่างระหว่างอิเล็กตรอนแต่ละตัวนานพอที่ไม่มีทางที่อิเล็กตรอนสองตัวจะปฏิสัมพันธ์กันได้ ผลการทดลองแสดงให้เห็นว่า ในช่วงแรก จุดสว่างจะปรากฏขึ้นบนฉากอย่างสุ่มเหมือนอนุภาคทั่วไป แต่เมื่อเวลาผ่านไปและอิเล็กตรอนสะสมจำนวนมากขึ้นเรื่อย ๆ (หลายหมื่นถึงหลายแสนตัว) ลายริ้วแทรกสอดควอนตัมที่ชัดเจนจะค่อย ๆ ปรากฏขึ้นมาอย่างเด่นชัด พิสูจน์ว่า \textbf{อิเล็กตรอนแต่ละตัวแทรกสอดกับตัวมันเองผ่านการเดินทางผ่านทั้งสองช่องทางพร้อม ๆ กันในสถานะซ้อนทับ}

\subsection{ความสมบูรณ์เสริมกันและการทำลายลายแทรกสอด}

หากผู้สังเกตการณ์พยายามติดตั้งตัวตรวจจับที่ช่องใดช่องหนึ่งเพื่อตรวจสอบว่าอิเล็กตรอนเดินทางผ่านช่องใดกันแน่ (การหาข้อมูล \textbf{Which-Way Information}) การวัดจะไปรบกวนสถานะควอนตัมและทำให้สถานะซ้อนทับสลายตัวไป เกิดปรากฏการณ์ที่เรียกว่า \textbf{การยุบตัวของฟังก์ชันคลื่น (Wavefunction Collapse)} ส่งผลให้พจน์การแทรกสอดกลายเป็นศูนย์ทันที ลายแทรกสอดจะอันตรธานหายไป และรูปแบบความน่าจะเป็นจะเปลี่ยนกลับไปเป็นผลบวกคลาสสิก $P_1(x) + P_2(x)$ ทันที สิ่งนี้สอดคล้องกับ \textbf{หลักความสมบูรณ์เสริมกัน (Complementarity Principle)} ของนีลส์ บอร์ ซึ่งระบุว่า สมบัติความเป็นคลื่น (การแทรกสอด) และสมบัติความเป็นอนุภาค (เส้นทางที่แน่นอน) เป็นสองแง่มุมที่ไม่สามารถสังเกตได้พร้อมกันในการทดลองเดียวกัน

% ==========================================================
% 2.7 พรมแดนงานวิจัยร่วมสมัย อินเทอร์เฟอโรเมทรีโมเลกุลชีวภาพ
% ==========================================================
\section{พรมแดนงานวิจัยร่วมสมัยและการทดลองระดับมหากาพย์}
\index{อินเทอร์เฟอโรเมทรีโมเลกุลขนาดใหญ่ (Macromolecule Interferometry)}
\index{โมเลกุลฟูลเลอรีน (Fullerene $C_{60}$)}
\index{ยาปฏิชีวนะกรามิซิดิน (Gramicidin)}
\index{การสูญเสียความพร้อมเพรียงเชิงควอนตัม (Decoherence)}
\index{การทดลองเลือกแบบหน่วงเวลา (Delayed-Choice Experiment)}

\subsection{การแทรกสอดควอนตัมของโมเลกุลขนาดใหญ่และโมเลกุลชีวภาพ}

คำถามสำคัญประการหนึ่งในฟิสิกส์พื้นฐานคือ ขอบเขตระหว่างโลกควอนตัม (ที่แสดงการซ้อนทับและการแทรกสอด) กับโลกมหภาคแบบคลาสสิก (ที่มีตำแหน่งแน่นอน) อยู่ที่ใด ในช่วงปี ค.ศ. 1999 ถึงปัจจุบัน กลุ่มวิจัยของ อันทอน ไซลิงเงอร์ (Anton Zeilinger, รางวัลโนเบลฟิสิกส์ 2022) และ มาร์คุส อาร์นดต์ (Markus Arndt) แห่งมหาวิทยาลัยเวียนนา ได้บุกเบิกการทดลองอินเทอร์เฟอโรเมทรีของคลื่นสสารในระดับโมเลกุลขนาดใหญ่ (Macromolecular Matter-Wave Interferometry)

\begin{enumerate}
    \item \textbf{การแทรกสอดของโมเลกุลฟูลเลอรีน ($C_{60}$ และ $C_{70}$)} ---- ในปี 1999 ประสบความสำเร็จในการสังเกตการแทรกสอดของโมเลกุลบัคกี้บอล $C_{60}$ ซึ่งมีมวลถึง $720~\text{amu}$ และประกอบด้วยอะตอมคาร์บอน 60 อะตอม
    \item \textbf{โมเลกุลสังเคราะห์ขนาดยักษ์ ($>25{,}000~\text{amu}$)} ---- ต่อมาได้ขยายการทดลองไปสู่โมเลกุลสารอินทรีย์ที่ออกแบบพิเศษซึ่งมีมวลมากกว่า $25{,}000~\text{Da}$ ประกอบด้วยอะตอมเกือบ 2,000 อะตอม
    \item \textbf{โมเลกุลชีวภาพธรรมชาติ กรามิซิดิน (Gramicidin)} ---- กลุ่มวิจัยของอาร์นดต์สามารถสาธิตการแทรกสอดควอนตัมของ \textbf{กรามิซิดิน} ซึ่งเป็นยาปฏิชีวนะเปปไทด์ธรรมชาติที่ประกอบด้วยกรดอะมิโน 15 ตัว มีมวลโมเลกุล $1882~\text{amu}$ โดยใช้เทคนิคการระเหิดด้วยเลเซอร์รังสีอัลตราไวโอเลตสั้นพิเศษ (Femtosecond Laser Desorption) ผสมกับลำอนุภาคก๊าซเฉื่อยเย็น และส่งผ่านเครื่องแทรกสอดคลื่นสสารแบบคาพิตซา-ดิแรก-ทัลบอต-เลา (Kapitza-Dirac-Talbot-Lau Interferometer)
\end{enumerate}

การทดลองเหล่านี้แสดงให้เห็นว่า แม้โมเลกุลจะมีขนาดใหญ่โต มีระดับขั้นความเสรีภายใน (การหมุนและการสั่น) นับร้อยรูปแบบ แต่หากสามารถป้องกันการรบกวนจากสิ่งแวดล้อมได้ จุดศูนย์กลางมวลของโมเลกุลชีวภาพก็ยังคงปฏิบัติตามกลศาสตร์ควอนตัมและแสดงการแทรกสอดของคลื่นเดอ บรอยล์ได้อย่างสมบูรณ์

\subsection{การสูญเสียความพร้อมเพรียงเชิงควอนตัมจากการแผ่รังสีความร้อน}

เหตุใดเราจึงไม่เห็นวัตถุในชีวิตประจำวันแสดงการแทรกสอดควอนตัม คำตอบทางฟิสิกส์ร่วมสมัยไม่ได้อยู่ที่การมีขีดจำกัดขนาดตายตัว หากแต่อยู่ที่ปรากฏการณ์ \textbf{การสูญเสียความพร้อมเพรียงเชิงควอนตัม (Quantum Decoherence)} 

เมื่อโมเลกุลมีอุณหภูมิสูง มันจะแผ่รังสีความร้อนวัตถุดำ (โฟตอนความร้อน) ออกมาในระหว่างที่เดินทาง โฟตอนแต่ละตัวที่ถูกเปล่งออกไปจะนำพาข้อมูลตำแหน่งของโมเลกุลรั่วไหลออกสู่สิ่งแวดล้อม เสมือนมีเครื่องวัดคอยแอบดูว่าโมเลกุลอยู่ที่ใด เมื่ออัตราการแผ่โฟตอนสูงขึ้น ความสัมพันธ์ทางเฟสระหว่างเส้นทางต่าง ๆ จะถูกทำลายลงอย่างรวดเร็ว ส่งผลให้คอนทราสต์ของลายแทรกสอดลดลงจนกลายเป็นศูนย์ การทำความเข้าใจกลไกเดโคฮีเรนซ์นี้เป็นกุญแจสำคัญที่สุดในการปกป้องสถานะควอนตัมในคอมพิวเตอร์ควอนตัมยุคปัจจุบัน

\subsection{การทดลองเลือกแบบหน่วงเวลาของวีลเลอร์และการลบควอนตัม}

ในคริสต์ศักราช 1978 จอห์น อาร์ชิบัลด์ วีลเลอร์ (John Archibald Wheeler, 1911--2008) ได้เสนอการทดลองในจินตนาการที่เรียกว่า \textbf{การทดลองเลือกแบบหน่วงเวลา (Delayed-Choice Experiment)} โดยเสนอว่า หากเราตัดสินใจว่าจะติดตั้งหรือถอดตัวตรวจจับเส้นทาง \textbf{หลังจาก} ที่โฟตอนหรืออิเล็กตรอนได้เดินทางผ่านช่องคู่ไปแล้ว อนุภาคจะรู้ได้อย่างไรว่ามันควรมีพฤติกรรมเป็นคลื่นหรืออนุภาค

การทดลองจริงได้รับการยืนยันสำเร็จในหลายห้องปฏิบัติการ รวมถึงการทดลองข้ามอวกาศโดยใช้ดาวเทียมสะท้อนแสงควอนตัม ซึ่งให้ผลตรงกันว่า \textbf{พฤติกรรมของอนุภาคในอดีตไม่ได้ถูกกำหนดไว้ล่วงหน้าอย่างตายตัว} จนกว่ากระบวนการวัดในขั้นตอนสุดท้ายจะเสร็จสมบูรณ์ และในการทดลอง \textbf{ยางลบควอนตัม (Quantum Eraser)} หากข้อมูลเส้นทางที่เคยถูกบันทึกไว้ถูกลบทำลายทิ้งไปโดยไม่สามารถกู้คืนได้ ลายแทรกสอดควอนตัมจะกลับคืนมาปรากฏขึ้นมาใหม่ได้อีกครั้ง

% ==========================================================
% 2.8 อินเทอร์เฟอโรเมทรีของนิวตรอนและอะตอมเย็น
% ==========================================================
\section{อินเทอร์เฟอโรเมทรีของนิวตรอนและอะตอมเย็นในสนามโน้มถ่วง}
\index{อินเทอร์เฟอโรเมทรีของนิวตรอน (Neutron Interferometry)}
\index{การทดลอง COW (COW Experiment)}
\index{อินเทอร์เฟอโรเมทรีของอะตอม (Atom Interferometry)}

\subsection{การทดลอง COW และการเลื่อนเฟสควอนตัมจากสนามโน้มถ่วง}

ในคริสต์ศักราช 1975 โรแบร์โต โคเลลลา (Roberto Colella) อัลเบิร์ต ดับเบิลยู. โอเวอร์เฮาเซอร์ (Albert W. Overhauser) และ แซมมูเอล เอ. เวอร์เนอร์ (Samuel A. Werner) ได้ทำการทดลองครั้งประวัติศาสตร์ที่รู้จักกันในนาม \textbf{การทดลอง COW (COW Experiment)} เพื่อทดสอบผลของสนามความโน้มถ่วงของโลกต่อฟังก์ชันคลื่นควอนตัมของอนุภาค

การทดลองใช้ \textbf{เครื่องแทรกสอดนิวตรอนผลึกซิลิกอนสมบูรณ์แบบ (Bonse-Hart Single-Crystal Silicon Interferometer)} ลำนิวตรอนความร้อนถูกแยกออกเป็นสองเส้นทางในแนวดิ่งที่ระดับความสูงต่างกัน $\Delta z$ ก่อนที่จะนำกลับมารวมกันเพื่อเกิดการแทรกสอด

พลังงานศักย์โน้มถ่วงของนิวตรอนมวล $m_n$ ที่ความสูง $z$ คือ $V(z) = m_n g z$ ส่งผลให้เวกเตอร์คลื่นของนิวตรอนในเส้นทางบนและเส้นทางล่างมีค่าไม่เท่ากัน ผลต่างเฟสเชิงควอนตัม ($\Delta\Phi_g$) ที่เกิดจากความโน้มถ่วงคำนวณได้จากการอินทิเกรตตามแนวเส้นทางปิด

\begin{equation}
\Delta\Phi_g = \frac{1}{\hbar}\oint \vec{p}\cdot d\vec{r} = -\frac{m_n^2 g \lambda A \sin\alpha}{\hbar^2}
\label{eq:cow_phase_shift}
\end{equation}

โดยที่ $A$ คือพื้นที่ที่ล้อมรอบด้วยเส้นทางเดินของลำนิวตรอน $\lambda$ คือความยาวคลื่นเดอ บรอยล์ของนิวตรอน และ $\alpha$ คือมุมเอียงของเครื่องแทรกสอดเทียบกับแนวระนาบ ผลการทดลองของ COW ยืนยันสมการ~\eqref{eq:cow_phase_shift} ด้วยความแม่นยำสูงมาก ซึ่งพิสูจน์ว่า \textbf{ศักย์ความโน้มถ่วงแบบดั้งเดิมของนิวตัน ($m g z$) สามารถทำให้เกิดการเลื่อนเฟสของฟังก์ชันคลื่นควอนตัมได้โดยตรง}

\subsection{อินเทอร์เฟอโรเมทรีของอะตอมเย็นและการวัดค่าคงตัวความโน้มถ่วง}

ในศตวรรษที่ 21 เทคโนโลยีการกักขังและทำให้อะตอมเย็นลงด้วยเลเซอร์ (Laser Cooling and Trapping) สู่ระดับไมโครเคลวินหรือนาโนเคลวิน (BEC) นำไปสู่การพัฒนา \textbf{อินเทอร์เฟอโรเมทรีของอะตอม (Atom Interferometry)} ซึ่งใช้พัลส์เลเซอร์แบบรามาน (Raman Transitions) ทำหน้าที่เป็นตัวแยกลำคลื่นและกระจกสะท้อนสำหรับคลื่นสสารของอะตอม (เช่น รูบิเดียม $^{87}\text{Rb}$ หรือซีเซียม $^{133}\text{Cs}$)

เทคนิคนี้มีความไวในการตรวจวัดความเร่งสูงกว่าเครื่องมือกลศาสตร์แบบเดิมหลายล้านเท่า ปัจจุบันถูกนำมาใช้เป็น \textbf{มาตรวัดความเร่งโน้มถ่วงควอนตัม (Quantum Gravimeter)} การทดสอบหลักความสมมูลของไอน์สไตน์ (Equivalence Principle) ในระดับความแม่นยำ $10^{-15}$ และการวัดค่าคงตัวโครงสร้างละเอียด (Fine-Structure Constant, $\alpha$) เพื่อทดสอบทฤษฎีอิเล็กโทรไดนามิกส์ควอนตัม (QED)

% ==========================================================
% 2.9 ทัศนศาสตร์อิเล็กตรอนขั้นสูงและเทคโนโลยี Cryo-EM
% ==========================================================
\section{ทัศนศาสตร์อิเล็กตรอนขั้นสูงและเทคโนโลยี Cryo-EM}
\index{กล้องจุลทรรศน์อิเล็กตรอนแบบส่องผ่าน (TEM)}
\index{กล้องจุลทรรศน์อิเล็กตรอนไครโอเจนิก (Cryo-EM)}
\index{ฟังก์ชันถ่ายโอนคอนทราสต์ (Contrast Transfer Function)}

\subsection{ขีดจำกัดการแยกขยายและฟังก์ชันถ่ายโอนคอนทราสต์ (CTF)}

ตามเกณฑ์ของเรย์ลีห์ (Rayleigh Criterion) กำลังแยกขยายเชิงทัศนศาสตร์ถูกจำกัดโดยการเลี้ยวเบนของคลื่น $\delta \approx 0.61 \frac{\lambda}{\text{NA}}$ สำหรับแสงที่ตามองเห็น ($\lambda \sim 500~\text{nm}$) กำลังแยกขยายจะถูกจำกัดอยู่ที่ประมาณ $200~\text{nm}$ ในขณะที่อิเล็กตรอนในกล้อง TEM ที่เร่งด้วยศักย์ $300~\text{kV}$ มีความยาวคลื่นเดอ บรอยล์เพียง $\lambda = 1.97~\text{pm} = 0.00197~\text{nm}$ ซึ่งเล็กกว่าแสงถึง 250,000 เท่า

อย่างไรก็ตาม ในทางปฏิบัติ กำลังแยกขยายของกล้องจุลทรรศน์อิเล็กตรอนไม่ได้ถูกจำกัดด้วยความยาวคลื่น แต่ถูกจำกัดด้วย \textbf{ความคลาดทางทัศนศาสตร์ (Aberrations)} โดยเฉพาะความคลาดทรงกลม (Spherical Aberration, $C_s$) ของเลนส์สนามแม่เหล็ก การส่งผ่านข้อมูลภาพอธิบายโดย \textbf{ฟังก์ชันถ่ายโอนคอนทราสต์ (Contrast Transfer Function --- CTF)}

\begin{equation}
\text{CTF}(k) = -\sin\left[ \frac{\pi}{2} C_s \lambda^3 k^4 - \pi \Delta f \lambda k^2 \right]
\label{eq:ctf_function}
\end{equation}

โดยที่ $k$ คือความถี่เชิงพื้นที่ $\Delta f$ คือค่าการปรับโฟกัสพร่อง (Defocus) และการปรับค่าให้เหมาะสมตามเงื่อนไขของเชอร์เซอร์ (Scherzer Defocus, $\Delta f_{\text{Sch}} = \sqrt{\frac{4}{3}C_s \lambda}$) ช่วยขยายย่านการส่งผ่านเฟสให้กว้างที่สุด

\subsection{การปฏิวัติชีววิทยาโครงสร้างด้วยกล้อง Cryo-EM}

การค้นพบระบบแก้ไขความคลาดทรงกลม ($C_s$-correctors) ร่วมกับเทคโนโลยีการแช่แข็งตัวอย่างทางชีวภาพอย่างรวดเร็ว (Vitrification) ในของเหลวอีเทนที่อุณหภูมิไนโตรเจนเหลว ($77~\text{K}$) นำไปสู่การพัฒนา \textbf{กล้องจุลทรรศน์อิเล็กตรอนไครโอเจนิก (Cryogenic Electron Microscopy --- Cryo-EM)} ซึ่งได้รับรางวัลโนเบลสาขาเคมีในปี ค.ศ. 2017 (Jacques Dubochet, Joachim Frank, และ Richard Henderson)

Cryo-EM สามารถถ่ายภาพอนุภาคเดี่ยว (Single-Particle Analysis) ของโมเลกุลชีวภาพที่ซับซ้อน เช่น ไวรัส ไรโบโซม และโปรตีนหนามของไวรัสโคโรนา (SARS-CoV-2 Spike Protein) จนได้โครงสร้างสามมิติที่มีความละเอียดระดับอะตอม ($< 1.2~\text{\AA}$) ซึ่งพิสูจน์ให้เห็นอย่างเป็นรูปธรรมว่า การควบคุมคลื่นสสารเดอ บรอยล์ของอิเล็กตรอนได้ปฏิวัติวงการวิทยาศาสตร์ชีวภาพและการแพทย์สมัยใหม่อย่างแท้จริง

% ==========================================================
% กล่องปฏิบัติการ WebAR
% ==========================================================
\begin{qrcodelabbox}{การทดลองเสมือนจริง การแทรกสอดช่องคู่ของอิเล็กตรอนเดี่ยว WebAR}{https://tsanaphysics.github.io/ar-Quantum/lab02_double_slit.html}
\begin{minipage}[c]{0.70\textwidth}
\textbf{คำแนะนำการทดลองเสมือนจริงสามมิติ (Interactive WebAR Simulation)}
\begin{enumerate}
    \item สแกน QR Code เพื่อเปิดแอปพลิเคชันจำลองการแทรกสอดช่องคู่แบบ 3 มิติ
    \item สังเกตการสะสมของจุดตกกระทบของอิเล็กตรอนเดี่ยวทีละตัวบนฉากรับ จนเกิดเป็นลายริ้วแทรกสอดควอนตัม
    \item เปิดใช้งานปุ่มสวิตช์ตรวจจับเส้นทาง (Which-Way Detector) เพื่อสังเกตการยุบตัวของฟังก์ชันคลื่นและการสลายตัวของลายแทรกสอด
    \item ปรับระยะห่างระหว่างช่องสลิต ($d$) และความยาวคลื่นเดอ บรอยล์ ($\lambda$) เพื่อสังเกตการขยายตัวของระยะห่างริ้วแทรกสอด
\end{enumerate}
\end{minipage}
\hfill
\begin{minipage}[c]{0.25\textwidth}
\centering
\qrcode[height=2.8cm]{https://tsanaphysics.github.io/ar-Quantum/lab02_double_slit.html}
\end{minipage}
\end{qrcodelabbox}

% ==========================================================
% ตัวอย่างโจทย์ P-S-C-T 5 ข้อเข้มข้น
% ==========================================================
\section{ตัวอย่างการแก้โจทย์เชิงวิเคราะห์ตามระเบียบวิธี P-S-C-T}

\begin{psctbox}{การคำนวณความยาวคลื่นเดอ บรอยล์เชิงสัมพัทธภาพในกล้อง TEM}
\textbf{Problem (โจทย์ปัญหา)}

กล้องจุลทรรศน์อิเล็กตรอนแบบส่องผ่านขั้นสูง (Transmission Electron Microscope --- TEM) ใช้ความต่างศักย์เร่งอิเล็กตรอน $V = 300~\text{kV}$ เพื่อให้ได้กำลังแยกขยายในระดับต่ำกว่าแองสตรอม
\begin{enumerate}
    \item จงคำนวณความยาวคลื่นเดอ บรอยล์ของอิเล็กตรอนโดยใช้สูตรกลศาสตร์ดั้งเดิม
    \item จงคำนวณความยาวคลื่นเดอ บรอยล์โดยใช้ทฤษฎีสัมพัทธภาพพิเศษอย่างถูกต้อง
    \item จงหาเปอร์เซ็นต์ความคลาดเคลื่อนหากละเลยผลทางสัมพัทธภาพ
\end{enumerate}

\textbf{Solution (วิธีทำ)}

\textbf{ขั้นตอนที่ 1} คำนวณแบบกลศาสตร์ดั้งเดิม (Classical)
พลังงานจลน์ $K = 300~\text{keV} = (300 \times 10^3)(1.6022 \times 10^{-19}~\text{J}) = 4.8066 \times 10^{-14}~\text{J}$
\begin{equation*}
p_{\text{class}} = \sqrt{2 m_e K} = \sqrt{2(9.1094 \times 10^{-31}~\text{kg})(4.8066 \times 10^{-14}~\text{J})} = 2.9592 \times 10^{-22}~\text{kg}\cdot\text{m/s}
\end{equation*}
\begin{equation*}
\lambda_{\text{class}} = \frac{h}{p_{\text{class}}} = \frac{6.62607 \times 10^{-34}}{2.9592 \times 10^{-22}} = 2.239 \times 10^{-12}~\text{m} = 2.239~\text{pm} = 0.02239~\text{\AA}
\end{equation*}

\textbf{ขั้นตอนที่ 2} คำนวณแบบสัมพัทธภาพพิเศษ (Relativistic)
พลังงานนิ่งของอิเล็กตรอนคือ $m_e c^2 = 510.999~\text{keV}$ พลังงานรวมคือ $E = K + m_e c^2 = 300 + 511.0 = 811.0~\text{keV}$
โมเมนตัมสัมพัทธภาพ $p_{\text{rel}} c = \sqrt{E^2 - (m_e c^2)^2} = \sqrt{(811.0)^2 - (511.0)^2} = \sqrt{657721 - 261121} = \sqrt{396600} \approx 629.76~\text{keV}$
\begin{equation*}
p_{\text{rel}} = \frac{629.76 \times 10^3 \times 1.6022 \times 10^{-19}~\text{J}}{2.9979 \times 10^8~\text{m/s}} = 3.3657 \times 10^{-22}~\text{kg}\cdot\text{m/s}
\end{equation*}
\begin{equation*}
\lambda_{\text{rel}} = \frac{h}{p_{\text{rel}}} = \frac{6.62607 \times 10^{-34}}{3.3657 \times 10^{-22}} = 1.9687 \times 10^{-12}~\text{m} = 1.969~\text{pm} = 0.01969~\text{\AA}
\end{equation*}

\textbf{ขั้นตอนที่ 3} คำนวณร้อยละความคลาดเคลื่อน
\begin{equation*}
\% \text{Error} = \frac{|\lambda_{\text{class}} - \lambda_{\text{rel}}|}{\lambda_{\text{rel}}} \times 100\% = \frac{|2.239 - 1.969|}{1.969} \times 100\% \approx 13.7\%
\end{equation*}

\textbf{Check (ตรวจสอบ)} ---- ความยาวคลื่นสัมพัทธภาพสั้นกว่าความยาวคลื่นคลาสสิกเสมอ เนื่องจากโมเมนตัมสัมพัทธภาพมีค่าสูงกว่า และความยาวคลื่นสั้นระดับ $1.97~\text{pm}$ เล็กกว่าขนาดอะตอมไฮโดรเจน ($53~\text{pm}$) อย่างมาก จึงให้กำลังแยกขยายระดับอะตอมได้

\textbf{Think (คิดต่อ)} ---- ความคลาดเคลื่อน $13.7\%$ ถือว่ามหาศาลมากสำหรับการศึกษาโครงสร้างผลึกด้วยเทคนิค Electron Diffraction หากวิศวกรไม่ใช้สูตรสัมพัทธภาพ การจำลองเลนส์แม่เหล็กไฟฟ้าในกล้อง TEM จะล้มเหลวโดยสิ้นเชิง
\end{psctbox}

\begin{psctbox}{การแทรกสอดของโมเลกุลยาปฏิชีวนะกรามิซิดิน}
\textbf{Problem (โจทย์ปัญหา)}

ในการทดลองของกลุ่มวิจัยอาร์นดต์ โมเลกุลยาปฏิชีวนะกรามิซิดิน ($C_{99}H_{140}N_{20}O_{17}$) มีมวลโมเลกุล $m = 1882~\text{amu} = 3.125 \times 10^{-24}~\text{kg}$ ถูกยิงด้วยความเร็วเฉลี่ย $v = 200~\text{m/s}$ ผ่านเกรตติงการเลี้ยวเบนที่มีระยะห่างช่อง $d = 266~\text{nm}$
\begin{enumerate}
    \item จงคำนวณความยาวคลื่นเดอ บรอยล์ของโมเลกุลกรามิซิดิน
    \item จงเปรียบเทียบความยาวคลื่นนี้กับขนาดทางกายภาพของโมเลกุล (ซึ่งมีขนาดยาวประมาณ $L_{\text{mol}} \approx 2.5~\text{nm}$)
    \item หากฉากรับอยู่ห่างจากเกรตติงเป็นระยะ $L = 1.20~\text{m}$ จงคำนวณระยะห่างระหว่างริ้วแทรกสอดบนฉาก ($\Delta y$)
\end{enumerate}

\textbf{Solution (วิธีทำ)}

\textbf{ขั้นตอนที่ 1} คำนวณโมเมนตัมและความยาวคลื่นเดอ บรอยล์
\begin{equation*}
p = m v = (3.125 \times 10^{-24}~\text{kg})(200~\text{m/s}) = 6.250 \times 10^{-22}~\text{kg}\cdot\text{m/s}
\end{equation*}
\begin{equation*}
\lambda = \frac{h}{p} = \frac{6.62607 \times 10^{-34}~\text{J}\cdot\text{s}}{6.250 \times 10^{-22}~\text{kg}\cdot\text{m/s}} = 1.060 \times 10^{-12}~\text{m} = 1.060~\text{pm} = 1060~\text{fm}
\end{equation*}

\textbf{ขั้นตอนที่ 2} เปรียบเทียบกับขนาดทางกายภาพของโมเลกุล
\begin{equation*}
\frac{\lambda}{L_{\text{mol}}} = \frac{1.060 \times 10^{-12}~\text{m}}{2.5 \times 10^{-9}~\text{m}} \approx 4.24 \times 10^{-4}
\end{equation*}
ความยาวคลื่นเดอ บรอยล์มีขนาดสั้นกว่าขนาดตัวของโมเลกุลเองถึงกว่า 2,300 เท่า

\textbf{ขั้นตอนที่ 3} คำนวณระยะห่างระหว่างริ้วแทรกสอด
\begin{equation*}
\Delta y = \frac{\lambda L}{d} = \frac{(1.060 \times 10^{-12}~\text{m})(1.20~\text{m})}{266 \times 10^{-9}~\text{m}} = 4.78 \times 10^{-6}~\text{m} = 4.78~\mu\text{m}
\end{equation*}

\textbf{Check (ตรวจสอบ)} ---- ระยะห่างริ้วแทรกสอด $4.78~\mu\text{m}$ มีขนาดใหญ่พอที่ระบบตรวจวัดรังสีไอออไนเซชันด้วยเลเซอร์ความละเอียดสูงสามารถสแกนตรวจนับได้อย่างแม่นยำ

\textbf{Think (คิดต่อ)} ---- แม้ความยาวคลื่นจะเล็กกว่าขนาดโมเลกุลมาก แต่การแทรกสอดเกิดขึ้นได้เพราะฟังก์ชันคลื่นของจุดศูนย์กลางมวลเกิดการแยกและซ้อนทับกันข้ามช่องเกรตติงที่มีระยะห่างถึง $266~\text{nm}$ ซึ่งกว้างกว่าขนาดโมเลกุลถึง 100 เท่า
\end{psctbox}

\begin{psctbox}{การแผ่กว้างออกของกลุ่มคลื่นอิเล็กตรอนเทียบกับอนุภาคฝุ่นมหภาค}
\textbf{Problem (โจทย์ปัญหา)}

พิจารณากลุ่มคลื่นเกาส์เซียนที่มีความไม่แน่นอนของตำแหน่งเริ่มต้น $\sigma_0 = 0.100~\text{nm}$ (ระดับขนาดอะตอม)
\begin{enumerate}
    \item จงคำนวณระยะเวลา ($t_d$) ที่ทำให้อนุภาคอิเล็กตรอนมีความกว้างของกลุ่มคลื่นเพิ่มขึ้นเป็นสองเท่า ($\sigma(t_d) = 2\sigma_0$)
    \item หากเปลี่ยนอนุภาคเป็นเม็ดฝุ่นขนาดจิ๋วที่มีมวล $m = 1.00 \times 10^{-9}~\text{kg}$ (หนึ่งไมโครกรัม) จงคำนวณเวลา $t_d$ เช่นเดียวกัน
\end{enumerate}

\textbf{Solution (วิธีทำ)}

\textbf{ขั้นตอนที่ 1} จากสมการการแผ่กว้างของกลุ่มคลื่น~\eqref{eq:wave_packet_spreading}
\begin{equation*}
\sigma(t) = \sigma_0 \sqrt{1 + \left( \frac{\hbar t}{2 m \sigma_0^2} \right)^2} = 2\sigma_0
\end{equation*}
ยกกำลังสองทั้งสองข้าง
\begin{equation*}
1 + \left( \frac{\hbar t_d}{2 m \sigma_0^2} \right)^2 = 4 \implies \frac{\hbar t_d}{2 m \sigma_0^2} = \sqrt{3} \implies t_d = \frac{2\sqrt{3} m \sigma_0^2}{\hbar}
\end{equation*}

\textbf{ขั้นตอนที่ 2} คำนวณสำหรับอิเล็กตรอน ($m_e = 9.109 \times 10^{-31}~\text{kg}$, $\sigma_0 = 10^{-10}~\text{m}$)
\begin{equation*}
t_d = \frac{2\sqrt{3}(9.109 \times 10^{-31}~\text{kg})(10^{-10}~\text{m})^2}{1.0546 \times 10^{-34}~\text{J}\cdot\text{s}} = \frac{3.1555 \times 10^{-50}}{1.0546 \times 10^{-34}} \approx 2.99 \times 10^{-16}~\text{วินาที}
\end{equation*}
กลุ่มคลื่นอิเล็กตรอนแผ่ขยายกว้างเป็นสองเท่าภายในเวลาเพียง $0.3$ เฟมโตวินาที

\textbf{ขั้นตอนที่ 3} คำนวณสำหรับเม็ดฝุ่น ($m = 1.00 \times 10^{-9}~\text{kg}$)
\begin{equation*}
t_d = \frac{2\sqrt{3}(1.00 \times 10^{-9}~\text{kg})(10^{-10}~\text{m})^2}{1.0546 \times 10^{-34}~\text{J}\cdot\text{s}} = \frac{3.464 \times 10^{-29}}{1.0546 \times 10^{-34}} \approx 3.28 \times 10^{5}~\text{วินาที} \approx 3.8~\text{วัน}
\end{equation*}

\textbf{Check (ตรวจสอบ)} ---- เวลาของอิเล็กตรอนสั้นอย่างยิ่ง สอดคล้องกับพฤติกรรมคลื่นที่แพร่กระจายรวดเร็ว ส่วนอนุภาคฝุ่นใช้เวลานานนับวัน

\textbf{Think (คิดต่อ)} ---- หากเม็ดฝุ่นมีขนาดระดับมิลลิกรัมหรือกรัม เวลา $t_d$ จะยาวนานกว่าอายุขัยของเอกภพ ($1.38 \times 10^{10}$ ปี) นี่คือเหตุผลทางฟิสิกส์ว่าเหตุใดเราจึงไม่เคยสังเกตเห็นวัตถุมหภาคแผ่กระจายตัวออกไปตามหลักกลศาสตร์ควอนตัม
\end{psctbox}

\begin{psctbox}{การเลี้ยวเบนของนิวตรอนเทอร์มอลและการทดสอบหลักความไม่แน่นอน}
\textbf{Problem (โจทย์ปัญหา)}

ลำนิวตรอนความร้อน (Thermal Neutrons) มวล $m_n = 1.675 \times 10^{-27}~\text{kg}$ มีพลังงานจลน์เฉลี่ยเท่ากับพลังงานความร้อน $K = \frac{3}{2}k_B T$ ที่อุณหภูมิห้อง $T = 300~\text{K}$ เคลื่อนที่ผ่านช่องเปิดเดี่ยว (Single Slit) ที่มีความกว้าง $a = 1.0~\mu\text{m}$
\begin{enumerate}
    \item จงคำนวณความเร็วเฉลี่ยและความยาวคลื่นเดอ บรอยล์ของนิวตรอน
    \item จงหามุมการเลี้ยวเบนของแถบมืดแรก ($\theta_1$) บนฉาก
    \item จงแสดงว่าผลคูณความไม่แน่นอนของตำแหน่งและโมเมนตัมที่ช่องเปิดสอดคล้องกับ $\Delta x \Delta p_y \sim h$
\end{enumerate}

\textbf{Solution (วิธีทำ)}

\textbf{ขั้นตอนที่ 1} คำนวณพลังงานจลน์และความยาวคลื่นเดอ บรอยล์
\begin{equation*}
K = \frac{3}{2}(1.3806 \times 10^{-23}~\text{J/K})(300~\text{K}) = 6.213 \times 10^{-21}~\text{J} = 0.0388~\text{eV}
\end{equation*}
\begin{equation*}
p_x = \sqrt{2 m_n K} = \sqrt{2(1.675 \times 10^{-27})(6.213 \times 10^{-21})} = 4.562 \times 10^{-24}~\text{kg}\cdot\text{m/s}
\end{equation*}
\begin{equation*}
\lambda = \frac{h}{p_x} = \frac{6.626 \times 10^{-34}}{4.562 \times 10^{-24}} = 1.452 \times 10^{-10}~\text{m} = 0.1452~\text{nm}
\end{equation*}

\textbf{ขั้นตอนที่ 2} หามุมการเลี้ยวเบนของแถบมืดแรกจาก $a \sin\theta_1 = \lambda$
\begin{equation*}
\sin\theta_1 \approx \theta_1 = \frac{\lambda}{a} = \frac{1.452 \times 10^{-10}~\text{m}}{1.0 \times 10^{-6}~\text{m}} = 1.452 \times 10^{-4}~\text{rad} \approx 0.00832^\circ
\end{equation*}

\textbf{ขั้นตอนที่ 3} ตรวจสอบหลักความไม่แน่นอน
ความไม่แน่นอนของตำแหน่งในแนวตั้งคือความกว้างของช่อง $\Delta y = a$
ความไม่แน่นอนของโมเมนตัมในแนวตั้งเกิดจากการกระจายมุมเลี้ยวเบน $\Delta p_y \approx p_x \sin\theta_1 \approx p_x \left(\frac{\lambda}{a}\right)$
\begin{equation*}
\Delta y \cdot \Delta p_y \approx a \cdot \left( p_x \frac{\lambda}{a} \right) = p_x \lambda = p_x \left( \frac{h}{p_x} \right) = h > \frac{\hbar}{2}
\end{equation*}

\textbf{Check (ตรวจสอบ)} ---- ผลคูณความไม่แน่นอนสอดคล้องกับหลักของไฮเซนเบิร์กอย่างสมบูรณ์

\textbf{Think (คิดต่อ)} ---- ยิ่งเราบีบช่องเปิดให้แคบลงเพื่อระบุตำแหน่ง $\Delta y$ ให้แม่นยำขึ้น ลำอนุภาคจะยิ่งเลี้ยวเบนกว้างออก ทำให้ $\Delta p_y$ กว้างขึ้นตาม เป็นเครื่องยืนยันว่าเราไม่สามารถกำจัดความไม่แน่นอนออกไปจากการทดลองได้
\end{psctbox}

\begin{psctbox}{อายุขัยของสถานะกระตุ้นและความกว้างธรรมชาติของเส้นสเปกตรัม}
\textbf{Problem (โจทย์ปัญหา)}

อะตอมโซเดียมถูกกระตุ้นขึ้นไปยังระดับพลังงาน $3p$ ซึ่งมีอายุขัยเฉลี่ย (Mean Lifetime) $\tau = 1.60 \times 10^{-8}~\text{s}$ (16.0 ns) ก่อนที่จะสลายตัวกลับสู่สถานะพื้น $3s$ โดยการปลดปล่อยโฟตอนแสงสีเหลืองที่มีความยาวคลื่น $\lambda_0 = 589.0~\text{nm}$
\begin{enumerate}
    \item จงคำนวณความไม่แน่นอนต่ำสุดของพลังงานของสถานะกระตุ้น ($\Delta E$) ในหน่วย eV
    \item จงคำนวณความกว้างธรรมชาติของเส้นสเปกตรัมในเทอมของความถี่ ($\Delta \nu$)
    \item จงคำนวณความกว้างธรรมชาติในเทอมของความยาวคลื่น ($\Delta \lambda$) และอัตราส่วน $\Delta \lambda / \lambda_0$
\end{enumerate}

\textbf{Solution (วิธีทำ)}

\textbf{ขั้นตอนที่ 1} คำนวณความไม่แน่นอนของพลังงานจาก $\Delta E \Delta t \ge \hbar/2$
โดยกำหนดให้ $\Delta t = \tau = 1.60 \times 10^{-8}~\text{s}$
\begin{equation*}
\Delta E \approx \frac{\hbar}{2\tau} = \frac{1.0546 \times 10^{-34}~\text{J}\cdot\text{s}}{2(1.60 \times 10^{-8}~\text{s})} = 3.296 \times 10^{-27}~\text{J} = \frac{3.296 \times 10^{-27}}{1.6022 \times 10^{-19}} \approx 2.057 \times 10^{-8}~\text{eV}
\end{equation*}

\textbf{ขั้นตอนที่ 2} แปลงความกว้างพลังงานเป็นความกว้างความถี่ $\Delta E = h \Delta \nu$
\begin{equation*}
\Delta \nu = \frac{\Delta E}{h} = \frac{\hbar / (2\tau)}{2\pi\hbar} = \frac{1}{4\pi\tau} = \frac{1}{4\pi(1.60 \times 10^{-8}~\text{s})} \approx 4.97 \times 10^6~\text{Hz} = 4.97~\text{MHz}
\end{equation*}

\textbf{ขั้นตอนที่ 3} คำนวณความกว้างในเทอมของความยาวคลื่น จาก $\nu = c/\lambda \implies |\Delta\lambda| \approx \frac{\lambda^2}{c}\Delta\nu$
\begin{equation*}
\Delta \lambda = \frac{(589.0 \times 10^{-9}~\text{m})^2}{2.998 \times 10^8~\text{m/s}}(4.97 \times 10^6~\text{s}^{-1}) \approx 5.75 \times 10^{-15}~\text{m} = 5.75 \times 10^{-6}~\text{nm}
\end{equation*}
\begin{equation*}
\frac{\Delta \lambda}{\lambda_0} = \frac{5.75 \times 10^{-15}~\text{m}}{589.0 \times 10^{-9}~\text{m}} \approx 9.76 \times 10^{-9}
\end{equation*}

\textbf{Check (ตรวจสอบ)} ---- ความกว้างธรรมชาติของเส้นสเปกตรัมแคบมาก (ประมาณหนึ่งในร้อยล้านส่วนของความยาวคลื่น) สอดคล้องกับเส้นสเปกตรัมที่คมชัดในอะตอมอิสระ

\textbf{Think (คิดต่อ)} ---- ในทางปฏิบัติ เส้นสเปกตรัมที่วัดได้ในห้องปฏิบัติการมักจะกว้างกว่านี้มาก เนื่องจากผลของการขยายตัวแบบดอปเปลอร์จากความร้อน (Doppler Broadening) และการชนกันของอะตอม (Collision Broadening)
\end{psctbox}

\begin{psctbox}{การทดลอง COW และการเลื่อนเฟสควอนตัมของนิวตรอนจากความโน้มถ่วง}
\textbf{Problem (โจทย์ปัญหา)}

ในเครื่องแทรกสอดนิวตรอนผลึกซิลิกอนของการทดลอง COW ลำนิวตรอนความร้อนมีความยาวคลื่นเดอ บรอยล์ $\lambda = 0.1445~\text{nm}$ เดินทางผ่านระนาบแทรกสอดรูปสี่เหลี่ยมขนมเปียกปูนที่มีพื้นที่ล้อมรอบ $A = 10.5~\text{cm}^2 = 1.05 \times 10^{-3}~\text{m}^2$
\begin{enumerate}
    \item จงคำนวณการเลื่อนเฟสเชิงควอนตัม ($\Delta\Phi_g$) เมื่อหมุนเอียงเครื่องแทรกสอดทำมุม $\alpha = 30.0^\circ$ เทียบกับแนวระนาบราบ
    \item จงคำนวณจำนวนริ้วแทรกสอดที่เลื่อนไป ($\Delta N = \Delta\Phi_g / 2\pi$)
\end{enumerate}

\textbf{Solution (วิธีทำ)}

\textbf{ขั้นตอนที่ 1} จากสมการการเลื่อนเฟสของ COW~\eqref{eq:cow_phase_shift}
\begin{equation*}
\Delta\Phi_g = \frac{m_n^2 g \lambda A \sin\alpha}{\hbar^2}
\end{equation*}

\textbf{ขั้นตอนที่ 2} แทนค่าคงตัวทางฟิสิกส์
\begin{align*}
m_n &= 1.67493 \times 10^{-27}~\text{kg} \\
g &= 9.80665~\text{m/s}^2 \\
\lambda &= 1.445 \times 10^{-10}~\text{m} \\
A &= 1.05 \times 10^{-3}~\text{m}^2 \\
\sin 30^\circ &= 0.5000 \\
\hbar &= 1.05457 \times 10^{-34}~\text{J}\cdot\text{s}
\end{align*}

คำนวณตัวเศษ
\begin{equation*}
\text{Numerator} = (1.67493 \times 10^{-27})^2 (9.80665)(1.445 \times 10^{-10})(1.05 \times 10^{-3})(0.5000) \approx 2.0913 \times 10^{-66}
\end{equation*}

คำนวณตัวส่วน
\begin{equation*}
\hbar^2 = (1.05457 \times 10^{-34})^2 \approx 1.1121 \times 10^{-68}
\end{equation*}

\textbf{ขั้นตอนที่ 3} คำนวณผลต่างเฟส
\begin{equation*}
\Delta\Phi_g = \frac{2.0913 \times 10^{-66}}{1.1121 \times 10^{-68}} \approx 188.04~\text{เรเดียน}
\end{equation*}

\textbf{ขั้นตอนที่ 4} คำนวณจำนวนริ้วที่เลื่อน
\begin{equation*}
\Delta N = \frac{\Delta\Phi_g}{2\pi} = \frac{188.04}{2(3.14159)} \approx 29.93~\text{ริ้ว} \approx 30~\text{ริ้ว}
\end{equation*}

\textbf{Check (ตรวจสอบ)} ---- การเลื่อนของริ้วประมาณ 30 ริ้วเป็นปริมาณมหภาคที่ตัวตรวจจับก๊าซ $^3\text{He}$ สามารถนับการแกว่งกวัดได้อย่างชัดเจน

\textbf{Think (คิดต่อ)} ---- การทดลองนี้แสดงให้เห็นว่าแม้ค่าคงตัวความโน้มถ่วง $g$ จะมีค่าน้อยมากเมื่อเทียบกับแรงทางนิวเคลียร์ แต่เนื่องจาก $\hbar$ มีขนาดเล็กระดับ $10^{-34}$ ผลคูณ $1/\hbar^2$ จึงขยายสัญญาณจนความโน้มถ่วงสามารถปรับเปลี่ยนพฤติกรรมของคลื่นสสารควอนตัมได้อย่างน่าอัศจรรย์
\end{psctbox}

\begin{psctbox}{กำลังแยกขยายของกล้อง Cryo-EM ภายใต้เงื่อนไขเชอร์เซอร์}
\textbf{Problem (โจทย์ปัญหา)}

กล้องจุลทรรศน์อิเล็กตรอนไครโอเจนิก (Cryo-EM) รุ่น Titan Krios ทำงานที่ความต่างศักย์ $V = 300~\text{kV}$ มีความยาวคลื่นอิเล็กตรอนเชิงสัมพัทธภาพ $\lambda = 1.97~\text{pm}$ และเลนส์ใกล้วัตถุมีค่าสัมประสิทธิ์ความคลาดทรงกลม $C_s = 2.70~\text{mm}$
\begin{enumerate}
    \item จงคำนวณหาค่าโฟกัสพร่องของเชอร์เซอร์ (Scherzer Defocus, $\Delta f_{\text{Sch}}$) เพื่อให้ได้คอนทราสต์เฟสที่ดีที่สุด
    \item จงประเมินกำลังแยกขยายเชิงจุดของเชอร์เซอร์ ($d_{\text{Sch}} = 0.66 C_s^{1/4} \lambda^{3/4}$)
\end{enumerate}

\textbf{Solution (วิธีทำ)}

\textbf{ขั้นตอนที่ 1} คำนวณค่าโฟกัสพร่องของเชอร์เซอร์
\begin{equation*}
\Delta f_{\text{Sch}} = \sqrt{\frac{4}{3} C_s \lambda} = \sqrt{\frac{4}{3}(2.70 \times 10^{-3}~\text{m})(1.97 \times 10^{-12}~\text{m})}
\end{equation*}
\begin{equation*}
\Delta f_{\text{Sch}} = \sqrt{7.092 \times 10^{-15}} \approx 8.42 \times 10^{-8}~\text{m} = 84.2~\text{nm}
\end{equation*}

\textbf{ขั้นตอนที่ 2} คำนวณกำลังแยกขยายเชิงจุด
\begin{align*}
C_s^{1/4} &= (2.70 \times 10^{-3})^{0.25} \approx 0.2279~\text{m}^{1/4} \\
\lambda^{3/4} &= (1.97 \times 10^{-12})^{0.75} \approx 1.656 \times 10^{-9}~\text{m}^{3/4} \\
d_{\text{Sch}} &= 0.66 \times (0.2279) \times (1.656 \times 10^{-9}) \approx 2.49 \times 10^{-10}~\text{m} = 0.249~\text{nm} = 2.49~\text{\AA}
\end{align*}

\textbf{Check (ตรวจสอบ)} ---- กำลังแยกขยาย $2.49~\text{\AA}$ อยู่ในระดับที่สามารถมองเห็นสายโซ่พอลิเปปไทด์ของโปรตีนได้อย่างชัดเจน

\textbf{Think (คิดต่อ)} ---- ในปัจจุบัน เมื่อติดตั้งระบบตัวแก้ความคลาดทรงกลมแบบเฮกซาโพล ($C_s$-corrector) ค่า $C_s$ จะลดลงเกือบเป็นศูนย์ ทำให้กำลังแยกขยายทะลุขีดจำกัดเชอร์เซอร์ลงไปถึงระดับต่ำกว่า $1.2~\text{\AA}$ ซึ่งสามารถระบุตำแหน่งอะตอมไฮโดรเจนเดี่ยวในโมเลกุลยาได้
\end{psctbox}

% ==========================================================
% 2.11 สรุปสาระสำคัญประจำบท
% ==========================================================
\section{สรุปสาระสำคัญประจำบทที่ 2}

\begin{enumerate}
    \item \textbf{สมมติฐานเดอ บรอยล์} ---- อนุภาคสสารทุกชนิดที่มีโมเมนตัม $p$ จะมีสมบัติคลื่นที่มีความยาวคลื่น $\lambda = h/p$ นำไปสู่ทวิภาคคลื่น-อนุภาคที่เป็นสมบัติสากลของธรรมชาติ
    \item \textbf{การยืนยันเชิงประจักษ์} ---- การทดลองของเดวิสสัน-เจอร์เมอร์ และจี.พี. ทอมสัน พิสูจน์การเลี้ยวเบนของอิเล็กตรอนในผลึกเดี่ยวและฟอยล์ทองคำ สอดคล้องกับเงื่อนไขของแบรกก์อย่างแม่นยำ
    \item \textbf{ความเร็วกลุ่มและความเร็วเฟส} ---- ความเร็วเฟสของคลื่นสสาร $v_p = c^2/v > c$ แต่ความเร็วกลุ่มของกลุ่มคลื่น $v_g = d\omega/dk = v$ ซึ่งเท่ากับความเร็วทางกายภาพของอนุภาคและไม่เกินความเร็วแสง
    \item \textbf{กลุ่มคลื่นเกาส์เซียน} ---- เป็นกลุ่มคลื่นที่มีผลคูณความไม่แน่นอนต่ำสุด $\Delta x \Delta p = \hbar/2$ แต่จะเกิดการกระจายตัวแผ่กว้างออกตามเวลา $\sigma(t) = \sigma_0\sqrt{1 + (\hbar t / 2m\sigma_0^2)^2}$
    \item \textbf{หลักความไม่แน่นอนของไฮเซนเบิร์ก} ---- เกิดจากการไม่สลับที่กันของตัวดำเนินการ $[\hat{x}, \hat{p}] = i\hbar$ กำหนดขีดจำกัดมูลฐานว่าไม่สามารถวัดตำแหน่งและโมเมนตัมพร้อมกันได้แม่นยำเกินกว่า $\Delta x \Delta p \ge \hbar/2$
    \item \textbf{การแทรกสอดของอนุภาคเดี่ยว} ---- อิเล็กตรอนเดี่ยวแทรกสอดกับตัวเองผ่านทั้งสองช่องทาง แต่หากตรวจวัดข้อมูลเส้นทาง (Which-Way Information) ฟังก์ชันคลื่นจะยุบตัวและลายแทรกสอดจะถูกทำลายทันที
    \item \textbf{พรมแดนโมเลกุลขนาดใหญ่} ---- การแทรกสอดของฟูลเลอรีน $C_{60}$ และโมเลกุลชีวภาพกรามิซิดินแสดงว่ากลศาสตร์ควอนตัมใช้ได้กับโมเลกุลขนาดยักษ์ โดยการเปลี่ยนผ่านสู่โลกคลาสสิกอธิบายได้ด้วยทฤษฎีการสูญเสียความพร้อมเพรียงเชิงควอนตัม (Decoherence)
\end{enumerate}

% ==========================================================
% คำถามเชิงมโนทัศน์ 20 ข้อ
% ==========================================================
\section{คำถามเชิงมโนทัศน์ประจำบทที่ 2}

\begin{enumerate}
    \item เหตุใดเราจึงไม่สังเกตเห็นการเลี้ยวเบนและการแทรกสอดของลูกฟุตบอลที่กำลังเคลื่อนที่ในชีวิตประจำวัน
    \item หากความเร็วของอนุภาคเข้าใกล้ความเร็วแสง ความยาวคลื่นเดอ บรอยล์จะเปลี่ยนแปลงไปอย่างไรเมื่อเทียบกับการคำนวณแบบคลาสสิก
    \item จงอธิบายว่าเหตุใดความเร็วเฟสของคลื่นเดอ บรอยล์จึงสามารถมีค่ามากกว่าความเร็วแสงได้โดยไม่ละเมิดทฤษฎีสัมพัทธภาพพิเศษ
    \item ในการทดลองเดวิสสัน-เจอร์เมอร์ หากให้ความร้อนแก่ผลึกนิกเกิลจนหลอมละลายกลายเป็นของเหลว ลายเลี้ยวเบนของอิเล็กตรอนจะยังคงปรากฏอยู่หรือไม่ เพราะเหตุใด
    \item จงอธิบายความแตกต่างเชิงแนวคิดระหว่าง การค้นพบของ เจ.เจ. ทอมสัน และการค้นพบของ จี.พี. ทอมสัน ผู้เป็นบุตรชาย
    \item เหตุใดกลุ่มคลื่นจึงจำเป็นต้องประกอบด้วยคลื่นหลายความถี่มารวมกัน และจะเกิดอะไรขึ้นกับความกว้างในปริภูมิโมเมนตัมหากกลุ่มคลื่นมีความกว้างในปริภูมิตำแหน่งแคบมาก
    \item อะไรคือสาเหตุทางกายภาพที่ทำให้กลุ่มคลื่นเกาส์เซียนแผ่กว้างออกเมื่อเวลาผ่านไป
    \item จงอธิบายว่าเหตุใดหลักความไม่แน่นอนของไฮเซนเบิร์กจึงไม่ใช่ความบกพร่องของเครื่องมือวัด แต่เป็นคุณลักษณะมูลฐานของธรรมชาติ
    \item อนุภาคสามารถหยุดนิ่งอย่างสมบูรณ์ที่ศูนย์สัมบูรณ์ ($T = 0~\text{K}$) ได้หรือไม่ จงอธิบายโดยใช้หลักความไม่แน่นอนของไฮเซนเบิร์ก
    \item จงอธิบายความสัมพันธ์ความไม่แน่นอนระหว่างพลังงานและเวลา และบทบาทในการเกิดอนุภาคเสมือนในสุญญากาศควอนตัม
    \item ในการทดลองช่องคู่ของโทโนมูระ หากยิงอิเล็กตรอนทีละตัวห่างกัน 1 ชั่วโมง จะเกิดลายแทรกสอดบนฉากรับหลังจากสะสมเป็นเวลานานหรือไม่ เพราะเหตุใด
    \item คำว่า การยุบตัวของฟังก์ชันคลื่น (Wavefunction Collapse) มีความหมายทางกายภาพอย่างไรในบริบทของการทดลองช่องคู่
    \item จงอธิบายหลักความสมบูรณ์เสริมกัน (Complementarity Principle) ของนีลส์ บอร์ โดยยกตัวอย่างการทดลองช่องคู่ประกอบ
    \item ในการทดลองของอาร์นดต์กับโมเลกุลกรามิซิดิน เหตุใดจึงต้องทำให้ห้องทดลองเป็นสุญญากาศสูงมากและต้องควบคุมอุณหภูมิไม่ให้โมเลกุลร้อนเกินไป
    \item การแผ่รังสีความร้อนของโมเลกุลขนาดใหญ่ทำลายลายแทรกสอดควอนตัมได้อย่างไรในเชิงทฤษฎีเดโคฮีเรนซ์
    \item จงอธิบายแนวคิดหลักของการทดลองเลือกแบบหน่วงเวลาของวีลเลอร์ (Wheeler's Delayed-Choice Experiment)
    \item ยางลบควอนตัม (Quantum Eraser) สามารถกู้คืนลายแทรกสอดที่หายไปได้อย่างไร
    \item ในกล้องจุลทรรศน์อิเล็กตรอนแบบส่องผ่าน (TEM) เหตุใดการใช้อิเล็กตรอนจึงให้กำลังแยกขยายสูงกว่าการใช้แสงอัลตราไวโอเลตหลายพันเท่า
    \item หากต้องการลดการกระจายตัวของกลุ่มคลื่นให้น้อยที่สุด เราควรออกแบบกลุ่มคลื่นเริ่มต้นอย่างไร
    \item สปินของอนุภาคมีผลต่อพฤติกรรมการแทรกสอดเชิงควอนตัมในการทดลองช่องคู่หรือไม่ อย่างไร
\end{enumerate}

% ==========================================================
% แบบฝึกหัดคำนวณและพิสูจน์เชิงวิเคราะห์ 20 ข้อ
% ==========================================================
\section{แบบฝึกหัดคำนวณและพิสูจน์เชิงวิเคราะห์ประจำบทที่ 2}

\begin{enumerate}
    \item จงคำนวณความยาวคลื่นเดอ บรอยล์ของโปรตอน ($m_p = 1.673 \times 10^{-27}~\text{kg}$) ที่มีพลังงานจลน์ $10.0~\text{MeV}$ โดยใช้ทฤษฎีสัมพัทธภาพพิเศษ
    \item ลำอิเล็กตรอนถูกเร่งผ่านความต่างศักย์ $V$ ตกกระทบผลึกที่มีระยะห่างระนาบ $d = 0.150~\text{nm}$ พบการเลี้ยวเบนอันดับที่หนึ่งที่มุม $\theta = 30.0^\circ$ จงหาค่าความต่างศักย์ $V$
    \item จงพิสูจน์ว่าสำหรับอนุภาคอิสระแบบสัมพัทธภาพ ผลคูณของความเร็วเฟสและความเร็วกลุ่มมีค่าเท่ากับ $c^2$ เสมอ ($v_p v_g = c^2$)
    \item ลูกกอล์ฟมวล $45.0~\text{g}$ เคลื่อนที่ด้วยความเร็ว $20.0~\text{m/s}$ วัดความเร็วได้แม่นยำ $0.01\%$ จงคำนวณหาความไม่แน่นอนต่ำสุดของตำแหน่งลูกกอล์ฟ
    \item อิเล็กตรอนถูกกักขังอยู่ในบริเวณขนาดอะตอมไฮโดรเจน ($\Delta x = 0.053~\text{nm}$) จงคำนวณหาความไม่แน่นอนต่ำสุดของความเร็ว และประเมินพลังงานจลน์ต่ำสุดของอิเล็กตรอน
    \item จงพิสูจน์ว่ากลุ่มคลื่นเกาส์เซียนตามสมการ~\eqref{eq:gaussian_packet_t0} สอดคล้องกับเงื่อนไขการนอร์แมลไลซ์ $\int_{-\infty}^\infty |\psi(x, 0)|^2 dx = 1$
    \item จงแสดงว่าการแปลงฟูเรียร์ของฟังก์ชันเกาส์เซียนตามสมการ~\eqref{eq:momentum_distribution} ให้ฟังก์ชันเกาส์เซียนในปริภูมิโมเมนตัม
    \item ในการทดลองช่องคู่ อิเล็กตรอนพลังงาน $100~\text{eV}$ ผ่านช่องคู่ที่มีระยะห่าง $d = 2.0~\mu\text{m}$ ไปตกกระทบบนฉากที่อยู่ห่างออกไป $L = 1.5~\text{m}$ จงคำนวณระยะห่างระหว่างริ้วแทรกสอดสว่างที่ติดกัน
    \item โมเลกุลฟูลเลอรีน $C_{60}$ ($m = 1.20 \times 10^{-24}~\text{kg}$) เคลื่อนที่ด้วยความเร็ว $v = 120~\text{m/s}$ ผ่านเกรตติงที่มีระยะห่าง $d = 100~\text{nm}$ จงคำนวณมุมเลี้ยวเบนอันดับที่หนึ่ง
    \item นิวตรอนพลังงานต่ำตัวหนึ่งมีความไม่แน่นอนของตำแหน่งเริ่มต้น $\sigma_0 = 1.00~\mu\text{m}$ จงคำนวณเวลาที่กลุ่มคลื่นนิวตรอนจะแผ่กว้างออกเป็น $10.0~\mu\text{m}$
    \item จงพิสูจน์ว่าค่าคาดหวังของโมเมนตัมสำหรับกลุ่มคลื่นเกาส์เซียนตามสมการ~\eqref{eq:gaussian_packet_t0} มีค่าเท่ากับ $\langle p \rangle = \hbar k_0$ และจงคำนวณหา $\langle p^2 \rangle$
    \item อนุภาค $\pi^0$ เมซอนมีมวลนิ่ง $135.0~\text{MeV}/c^2$ และมีอายุขัยเฉลี่ย $\tau = 8.4 \times 10^{-17}~\text{s}$ จงคำนวณหาความไม่แน่นอนของพลังงานนิ่ง และอัตราส่วน $\Delta E / (m c^2)$
    \item ในการทดลองของมาร์คุส อาร์นดต์ หากโมเลกุลมีอุณหภูมิภายใน $T = 1000~\text{K}$ และปลดปล่อยโฟตอนความร้อน 1 ตัวความยาวคลื่น $\lambda_{\text{thermal}} = 5.0~\mu\text{m}$ จงคำนวณโมเมนตัมสะท้อนที่ถ่ายโอนให้แก่โมเลกุล
    \item ในกล้อง TEM ความต่างศักย์เร่ง $200~\text{kV}$ จงคำนวณความเร็วของอิเล็กตรอนในรูปของสัดส่วนความเร็วแสง ($v/c$) และความยาวคลื่นเดอ บรอยล์
    \item จงแสดงว่าสำหรับตัวดำเนินการแฮร์มิเชียน $\hat{A}$ และ $\hat{B}$ หาก $[\hat{A}, \hat{B}] = 0$ จะไม่มีข้อจำกัดความไม่แน่นอนตามหลักของไฮเซนเบิร์ก ($\Delta A \Delta B \ge 0$)
    \item ลำอะตอมฮีเลียมที่อุณหภูมิ $T = 2.0~\text{K}$ เคลื่อนที่ออกจากหัวฉีด จงคำนวณความยาวคลื่นเดอ บรอยล์เฉลี่ยของอะตอมฮีเลียม
    \item จงคำนวณหาค่าความต่างศักย์เร่งที่ทำให้อิเล็กตรอนมีความยาวคลื่นเดอ บรอยล์เท่ากับความยาวคลื่นคอมป์ตันของมันพอดี ($\lambda = \lambda_C$)
    \item อนุภาคกล่องหนึ่งมิติความยาว $L$ มีฟังก์ชันคลื่นสถานะพื้น $\psi(x) = \sqrt{\frac{2}{L}}\sin\left(\frac{\pi x}{L}\right)$ จงคำนวณหา $\Delta x$ และ $\Delta p$ และตรวจสอบว่าสอดคล้องกับหลักความไม่แน่นอนหรือไม่
    \item จงคำนวณหาแรงดันความร้อนและผลของการเลื่อนดอปเปลอร์ในอินเทอร์เฟอโรเมทรีของโมเลกุลชีวภาพ
    \item จงแสดงว่าความเร็วกลุ่มสำหรับคลื่นในตัวกลางที่มีการกระจายตัว (Dispersive Medium) สัมพันธ์กับดรรชนีหักเห $n(\lambda)$ ตามสมการ $v_g = \frac{c}{n - \lambda (dn/d\lambda)}$
\end{enumerate}

% ==========================================================
% โครงงานจำลองเชิงคำนวณด้วย Python
% ==========================================================
\section{โครงงานจำลองเชิงคำนวณด้วยภาษา Python}

โปรแกรมภาษา Python ต่อไปนี้ใช้จำลองการวิวัฒนาการเชิงเวลาและการแผ่กว้างออกของกลุ่มคลื่นเกาส์เซียนเชิงควอนตัม (Wave Packet Dispersion) ตามสมการ~\eqref{eq:wave_packet_time_evolved}

\begin{lstlisting}[language=Python, caption={โปรแกรม Python จำลองการแผ่กระจายตัวของกลุ่มคลื่นเกาส์เซียนควอนตัมเชิงเวลา}]
import numpy as np
import matplotlib.pyplot as plt

# กำหนดหน่วยทางฟิสิกส์ธรรมชาติ (Atomic / Normalized units)
hbar = 1.0     # Reduced Planck constant
m = 1.0        # Particle mass
sigma_0 = 1.0  # Initial spatial width (sigma_0)
k_0 = 5.0      # Initial wavevector (momentum p_0 = 5.0)

# กำหนดอาณาเขตของพิกัด x
x = np.linspace(-15, 35, 1000)

def wave_packet(x, t, sigma_0, k_0, m, hbar):
    """คำนวณฟังก์ชันคลื่นและความหนาแน่นความน่าจะเป็น ณ เวลา t"""
    alpha = hbar * t / (2.0 * m * sigma_0**2)
    denom = 1.0 + 1j * alpha
    v_0 = hbar * k_0 / m
    
    # คำนวณฟังก์ชันคลื่นไซ (Psi)
    psi = (1.0 / (2.0 * np.pi * sigma_0**2))**0.25 * (1.0 / np.sqrt(denom)) * \
          np.exp(- (x - v_0 * t)**2 / (4.0 * sigma_0**2 * denom)) * \
          np.exp(1j * (k_0 * x - (hbar * k_0**2 / (2.0 * m)) * t))
    
    prob_density = np.abs(psi)**2
    sigma_t = sigma_0 * np.sqrt(1.0 + alpha**2)
    return prob_density, v_0 * t, sigma_t

# จำลองที่เวลาต่าง ๆ กัน
times = [0.0, 1.5, 3.0, 4.5]
colors = ['#00F0FF', '#7B2CBF', '#FF007F', '#00FF66']

plt.figure(figsize=(10, 6))

for t, col in zip(times, colors):
    prob, x_center, sig_t = wave_packet(x, t, sigma_0, k_0, m, hbar)
    plt.plot(x, prob, label=f't = {t:.1f} (center = {x_center:.1f}, width = {sig_t:.2f})', color=col, lw=2.2)
    plt.fill_between(x, prob, alpha=0.15, color=col)

plt.title('Time Evolution and Dispersion of Gaussian Wave Packet', fontsize=14, fontweight='bold')
plt.xlabel('Position x (arbitrary units)', fontsize=12)
plt.ylabel('Probability Density $|\\psi(x, t)|^2$', fontsize=12)
plt.grid(True, linestyle=':', alpha=0.6)
plt.legend(fontsize=10)
plt.tight_layout()
plt.savefig('gaussian_wavepacket_dispersion.png', dpi=300)
plt.show()
\end{lstlisting}
